\documentclass[11pt]{article}
\usepackage[a4paper,margin=27mm]{geometry}
\usepackage{amsmath,amssymb,amsthm,hyperref}
\newtheorem{theorem}{Theorem}
\newtheorem{lemma}[theorem]{Lemma}
\newtheorem{corollary}[theorem]{Corollary}
\newcommand{\tp}{\nu_{\triangle}}
\newcommand{\tc}{\tau_{\triangle}}
\newcommand{\mc}{\operatorname{MaxCut}}
\title{Two neighborhood types in Tuza's conjecture\\
\large Direct covers, quantitative margins, and compatible packings}
\author{Working research notes\\OpenAI Codex (GPT-6 Astra)}
\date{8 October 2026}
\begin{document}
\maketitle

\noindent\textbf{Status (8 October 2026).} Theorem~\ref{thm:two-types-all}
is fully formalized and verified in Lean 4.34.0: every finite simple split
graph with at most two active independent-side neighborhood types satisfies
Tuza's inequality, with no restriction on clique size, overlap, or multiplicities.
The final route uses matching allocations, explicit cut covers, greedy core
completion via Mantel's bound, and sum-coloring of triangles.
A finite certificate covers all $23\,848\,371$ required compressed parameter
tuples for $2\leq k\leq42$. Analytic reductions and a continuous certificate
with $178\,506$ leaves cover every $k\geq43$ simultaneously.
The final graph theorem and the certificates, including their soundness and
coverage proofs, are kernel-checked. Only the standard axioms
\texttt{propext}, \texttt{Classical.choice}, and \texttt{Quot.sound} occur.
The final formal-verification section records the evidence.

These are development notes: intermediate sections preserve earlier routes,
stronger bounds, failed comparisons and historical verification checkpoints.
Their partial-status comments describe those checkpoints, not the completed
formalization. The final route does not use the exact clique-packing formula,
a threshold-graph theorem, the exact cover formula, or permutation averaging.
This proves the stated two-type subclass, not Tuza's conjecture for arbitrary
graphs. No publication priority or independent expert review is claimed.

\section{The graph class and notation}
All graphs are finite, simple, and undirected. Let $G$ have a split partition
$V(G)=C\mathbin{\dot\cup}I$, with $C$ a clique of size $k$ and $I$ independent.
After deleting vertices of $I$ with degree at most one, assume that every
remaining vertex has neighborhood $S$ or $T$ in $C$. Write $p,q$ for their
multiplicities and
\[
 a=|S\setminus T|,\quad b=|T\setminus S|,\quad
 u=|S\cap T|,\quad o=|C\setminus(S\cup T)|.
\]
Thus $k=a+b+u+o$, $s=|S|=a+u$, and $t=|T|=b+u$.
The target is $\tc(G)\leq 2\tp(G)$ for all six parameters.
Here $\tc$ is the minimum size of an edge set meeting every triangle, and
$\tp$ is the maximum number of pairwise edge-disjoint triangles.

Nested neighborhoods give threshold graphs, already covered by
Bonamy et al.~\cite{bonamy}. The cases beyond that known result have $p,q>0$ and $a,b>0$.
Zeng~\cite{zeng} treats two types with a specified eight-vertex clique part.
The purpose here is to develop reductions with no bound on $k$.

\section{The final proof route}
Remove outside vertices that participate in no triangle, and compress each
remaining multiplicity with Theorem~\ref{thm:compression}. The final
proof has only the following branches:
\[
 \begin{array}{ccl}
 k\leq1&:&\text{there are no triangles},\\
 2\leq k\leq42&:&\text{the sum-coloring finite certificate},\\
 k\geq43&:&\text{analytic reductions and the unchanged continuous certificate}.
 \end{array}
\]
Lemma~\ref{lem:cyclic-reductions} handles one-type graphs, a small
neighborhood, small total multiplicity, and two neighborhoods no larger
than half the clique. Every remaining large-core graph lies in the
domain of Lemma~\ref{lem:cyclic-interval}.

The certificates compare explicit covers (balanced cuts, shadow cuts,
and the unbalanced common cut) with packing lower bounds from local
matching colors, greedy core completion, and sum-coloring of triangles.
Their comparison gives $\tc\leq2\tp$. The matching decomposition,
Mantel bound, cut rounding, and sum-coloring arguments are given in the notes.

The exact clique-packing formula, threshold-graph theorem, exact cover
formula, permutation averaging, extremal theorem for two holes, moment calculations, stability
and lifting arguments, multiplicity balancing, and separate equal-size
theorem are not dependencies of this final route. They remain in the
notes as background or earlier research. The elementary cyclic bound is
weaker than the exact clique optimum; its use is justified by new
analytic constants and fresh certificates, not by reusing earlier
numerical conclusions. The subsequent sum-coloring simplification leaves
the continuous bounds unchanged and has its own new finite check.

\section{A sharp linear bound on twin multiplicities}
The next result applies to a more general graph. Let $H$ be any finite
simple graph and let $S\subseteq V(H)$ have size $s\geq2$. Form $G_m$ by
adding $m$ independent vertices, each with neighborhood exactly $S$.
The set $S$ need not be a clique.
Define
\[
 h(s)=\chi'(K_s)=
 \begin{cases}s-1,&s\text{ even},\\s,&s\text{ odd}.\end{cases}
\]
For completeness, the required round-robin decomposition is explicit.
If $s$ is even, label the vertices by $\{\infty\}\cup\mathbb Z/(s-1)\mathbb Z$.
For each residue $c$, take the matching
\[
 \{\{\infty,c\}\}\ \cup\
 \{\{c+i,c-i\}:1\leq i\leq(s-2)/2\}.
\]
It is perfect: the nonzero residues modulo the odd number $s-1$ occur
exactly once in pairs $\{i,-i\}$. Every edge not incident with $\infty$
belongs to exactly one matching, because its endpoints $a,b$ determine
$c=(a+b)/2$ modulo $s-1$; division by $2$ is valid for an odd modulus.
Edges incident with $\infty$ have their evident unique color. This also
covers $s=2$, when the second set is empty. For odd $s$, use this
construction on $s+1$ vertices and delete one vertex. The resulting
$s$ matchings each have $(s-1)/2$ edges. Thus every color class has
$\lfloor s/2\rfloor$ edges. The degree lower bound for even $s$, and
edge counting with at most $\lfloor s/2\rfloor$ edges per color for odd
$s$, also prove that the displayed number of colors is minimal.

There is a shorter route to the class-size property once properness is
known. In any proper coloring of $K_s$ with $h(s)$ available colors, each
class $M_c$ is a matching, so $|M_c|\leq\lfloor s/2\rfloor$. But
\[
 \sum_{c=1}^{h(s)}|M_c|=\binom{s}{2}
 =h(s)\lfloor s/2\rfloor.
\]
Thus every class attains its upper bound. This counting argument avoids
counting the individual residue classes. It is now verified in Lean for
every proper coloring with this palette size, not just the particular
round-robin construction. Consequently every set of $r$ colors contains
exactly $r\lfloor s/2\rfloor$ edges.

\begin{corollary}[Allocation on one clique neighborhood]\label{cor:one-neighborhood-allocation}
Suppose $S$ is a clique in $H$, with $s=|S|\geq2$. For every $m\geq0$,
the graph $G_m$ has an edge-disjoint triangle packing of exactly
\[
 \min(m,h(s))\lfloor s/2\rfloor
\]
triangles, each using one new vertex. In particular, this expression is
a lower bound for $\tp(G_m)$.
\end{corollary}
\begin{proof}
Put $r=\min(m,h(s))$ and choose any $r$ colors of a proper
$h(s)$-coloring of $K_S$. Every chosen class has $\lfloor s/2\rfloor$
edges. Assign the chosen colors to $r$ distinct new vertices and form a
triangle over each selected edge with its assigned center. The old edges
are distinct, triangles at one center use a matching, and triangles at
different centers have distinct spokes. All their edges are present
because $S$ is a clique. The count is $r\lfloor s/2\rfloor$, including
the empty packing when $m=0$.
\end{proof}
This asserts the size of a constructed packing, not the optimum over all
packings: triangles wholly inside $H$ can contribute additional members.
The formula cannot simply be added for two overlapping neighborhoods,
because their selected old edges may coincide.

\paragraph{Lean verification of uniform classes and allocation.}
The module \path{lean-matching-classes/MatchingClasses.lean} retains
the previous graph-packing source verbatim and proves the counting
argument above. Its theorem \texttt{uniform\_class\_size} applies to any
proper coloring of all two-element subsets of a finite neighborhood.
The theorem \texttt{selected\_edges\_card} proves the exact size for
every selected palette, without an averaging step.

The theorem \texttt{palette\_packing} assigns selected colors injectively
to actual new centers. It proves adjacency, pairwise edge-disjointness,
exact cardinality, and the presence of a new center in every produced
triangle. The theorem \texttt{clique\_neighborhood\_packing} then handles
every $m$, including $m=0$ and $m>h(s)$, in an arbitrary ambient vertex
type. Its finite-graph corollary
\texttt{clique\_neighborhood\_lower\_bound} compares that size to the
actual attained graph optimum. Unlike multiplicity compression, this
allocation theorem requires $S$ to be a clique.

Azure job \texttt{d209d4a414c24fc2} verified the complete source on its
first attempt in $72.224$ seconds, using Lean 4.34.0 and the pinned Mathlib
revision. All seven final dependency audits use only \texttt{propext},
\texttt{Classical.choice}, and \texttt{Quot.sound}; the arithmetic
palette identity does not need \texttt{Classical.choice}. The false
arithmetic control was rejected and the returned source matches the
local source. The worker had a 16 GiB memory limit, a 25-minute runtime
limit, and an independent cloud deletion guard. Cleanup evidence is
recorded separately from proof success.

\textbf{Assessment.} Uniform matching-class sizes and the entire
single-neighborhood allocation bridge are now kernel-checked. The
independent two-neighborhood allocation below now also has a verified
compatibility construction and overlap estimate. The coupled-color lower
bound below is now verified too, with coherent labels and finite cyclic
averaging. Both constructions control shared edges before combining
the two types.

\begin{theorem}[Exact multiplicity compression]\label{thm:compression}
For every $m\geq0$,
\[
 \tp(G_m)=\tp(G_{\min(m,h(s))}),\qquad
 \tc(G_m)=\tc(G_{\min(m,s-1)}).
\]
Consequently truncating at $h(s)$ preserves both parameters simultaneously.
\end{theorem}
\begin{proof}
Both parameters are nondecreasing under addition of vertices, so it suffices
to prove the reverse inequalities above their proposed thresholds.

Take any triangle packing in $G_m$. Every triangle using a new vertex uses
one edge of $H[S]$. Collect these edges into a graph $J\subseteq H[S]$.
They are all distinct. Restrict a fixed $h(s)$-edge-coloring of $K_S$ to $J$.
Replace the center of each such triangle by the new vertex indexed by the
color of its clique edge. Triangles assigned to the same center use a matching,
so their spokes are distinct. Triangles assigned to different centers have
different spokes, and their clique edges are distinct. All triangles lying
inside $H$ are left unchanged; they use none of these clique edges and no
new spokes. Thus the packing survives on $h(s)$ new vertices with the same
cardinality.

For covers, let $r=s-1$ and take any triangle cover $F$ of $G_r$.
Put $F_H=F\cap E(H)$ and let $J$ consist of the edges of $H[S]$ not in $F_H$.
For every new vertex $x$, the set $Z_x=\{v\in S:xv\in F\}$ is a vertex cover
of $J$: the triangle formed by $x$ and an edge of $J$ must be met by a
selected spoke. Choose $x_*$ minimizing $|Z_x|$. Every edge of $J$ is
incident to $Z_{x_*}$, and each vertex of $S$ has at most $s-1$ neighbors
in $S$. Consequently
\[
 |E(J)|\leq (s-1)|Z_{x_*}|
 \leq \sum_{x=1}^{r}|Z_x|.
\]
The sum counts distinct selected spokes. Replace them all by $E(J)$.
The resulting set $F_H\cup E(J)$ has size at most $|F|$, still covers every
triangle in $H$, and contains every edge of $H[S]$. It therefore covers the
triangles centered at arbitrarily many new vertices. This directly
transforms every cover, without introducing a separate vertex-cover optimum.
\end{proof}

\paragraph{A reusable recoloring statement.}
The packing part has a useful formulation that separates its construction
from the extraction of data from a graph. Let $S$ be an $s$-element subset
of an arbitrary old vertex set $V$, with $s\geq2$. Let $Q$ be an
edge-disjoint family of triples in $V$, and let $J\subseteq\binom S2$.
Assume no member of $J$ is an edge of a member of $Q$. A proper
$h(s)$-edge-coloring $c$ gives the explicit new family
\[
 P=Q\ \cup\ \{e\cup\{x_{c(e)}\}:e\in J\},
\]
where the $h(s)$ vertices $x_i$ are new and distinct. Then $P$ is
edge-disjoint and $|P|=|Q|+|J|$. All members of $Q$ are retained.
Indeed, two new triangles can share a spoke only if their core edges
meet and have the same color, which properness excludes for distinct
edges. A new triangle can share an edge with an old one only on its
core edge, excluded by the assumption on $J$. Every triangle is distinct,
so the cardinalities add. This formulation permits the triangles of the
second neighborhood type to remain in $Q$ while the first type is
recolored; it places no bound on $V$, $Q$, or the original multiplicity.

\paragraph{Lean verification of the recoloring construction.}
The module \path{lean-compression/Compression.lean} proves the proper
coloring bound for every finite size $s\geq2$, then proves the preceding
recoloring statement for actual triples in the disjoint union
$V\sqcup\operatorname{Fin}(h(s))$. The final theorem is
\texttt{recolor\_on\_neighborhood}. It includes the exact cardinality,
pairwise disjointness of two-element edge sets, retention of every old
triple, and a centered triangle over every selected core edge. No bound
on the ambient vertex set is assumed. The proof handles an arbitrary
finite neighborhood inside that ambient set, not just an isolated clique.

Azure job \texttt{8afa912918af43f0} compiled the complete module with Lean
4.34.0 and the pinned Mathlib revision in $68.471$ seconds. All four
audited results use only \texttt{propext}, \texttt{Classical.choice}, and
\texttt{Quot.sound}; the false arithmetic control was rejected. A prior
initial module passed separately; the first extended attempt failed on
two elaboration errors and is retained as a failure. Memory and time
limits and independent cloud cleanup were used; actual deletion status
is recorded separately in the run reports.

This module verifies the constructive packing component of
Theorem~\ref{thm:compression}. Its graph-level packing equality is now
verified by the companion module described below. A further module
verifies cover compression at $s-1$ centers and simultaneous preservation
of both parameters. Uniform matching-class sizes and single-neighborhood
allocation are now verified by the module described above. No scalar expression or numerical certificate
domain changed in these formalization steps.

\paragraph{Extracting the data from an arbitrary graph packing.}
Let $P$ be a triangle packing in $G_m$. Each triangle has at most one new
vertex, since the new vertices are independent. Its old part therefore
has at least two vertices. The old-part map is injective on $P$: equal
old parts would contain a shared two-element edge. Split $P$ into old
triangles and triangles using one new vertex. Their old parts give an
old packing $Q$ and a family $J\subseteq E(H)\cap\binom S2$, respectively,
with $|P|=|Q|+|J|$. No edge in $J$ lies in a triangle of $Q$, again by
edge-disjointness of $P$. Thus the hypotheses of the recoloring
construction follow from the input packing itself.

Every produced triple is a triangle of the target graph: its old edge
belongs to $H$, and both endpoints lie in $S$, so both new spokes are
present. When $m\geq h(s)$, the reverse map simply embeds the $h(s)$ new
vertices into the $m$ new vertices. Thus every attainable packing size
is preserved in both directions, and so is the maximum. When $m\leq h(s)$,
the truncated graph is unchanged. This packing argument works even if
$S$ is not a clique: only the actually selected edges of $H$ are used,
while the complete-graph coloring supplies a convenient common palette.

\paragraph{Lean verification of exact packing compression.}
The companion \path{lean-graph-compression/GraphCompression.lean}
formalizes the whole graph bridge. It defines $G_m$ as a Mathlib
\texttt{SimpleGraph}, proves that the extracted data satisfy every
recoloring hypothesis, and checks graph adjacency of the output triples.
An injective graph map gives the reverse direction. Its theorem
\texttt{packing\_sizes\_attached\_min} proves equivalence of every
attainable finite packing size, including when the old vertex type is
infinite. For finite old vertex types,
\texttt{packingNumber\_spec} proves that the finite supremum defining the
packing number is attained and bounds every packing. The final theorem
\texttt{packingNumber\_attached\_min} therefore states the exact equality
\[
 \tp(G_m)=\tp(G_{\min(m,h(s))})
 \qquad(s\geq2,\ m\geq0).
\]
It assumes neither that $S$ is a clique nor a bound on $m$ or the size
of the old graph. This is a general compression lemma; the full
Tuza inequality remains a separate target.

Azure job \texttt{b1553d10589540b3} verified the combined source with Lean
4.34.0 and the same pinned Mathlib revision in $71.784$ seconds. All six
audited graph-level results depend only on \texttt{propext},
\texttt{Classical.choice}, and \texttt{Quot.sound}. The false arithmetic
control was rejected and the returned source is byte-identical to the
local source. The first attempt failed on elaboration errors and is
retained as a failure; a local preparation superseded before launch
created no cloud resources. Every actual job had bounded memory and time
and an independent cloud deletion guard; deletion evidence is recorded
separately from proof success.

\paragraph{Lean verification of exact cover compression.}
The module \path{lean-cover-compression/CoverCompression.lean} retains
the preceding graph-packing source verbatim and adds the cover proof.
Its definition \texttt{IsCover} requires every selected pair to be an
actual graph edge and every triangle to contain a selected edge.
The lemma \texttt{replace\_spokes} implements the minimum-selected-spoke
argument above. The old-edge cover it produces contains every edge of
$H[S]$, so it extends to any number of new centers at no extra cost.
Restriction along an injective graph map supplies the converse direction.

The definition \texttt{coverNumber} is the least feasible edge budget;
\texttt{coverNumber\_spec} proves that this value is attained and is no
greater than the size of any valid cover. The theorem
\texttt{coverNumber\_attached\_min} therefore proves the exact equality
of triangle-cover optima at the threshold $s-1$. The final theorem
\texttt{both\_parameters\_attached\_min} combines both halves of
Theorem~\ref{thm:compression} at the common threshold $h(s)$.
The result quantifies over every finite old graph, every neighborhood
of size at least two, and every multiplicity.

Azure job \texttt{8e4f6877a2034648} verified the complete combined source
with Lean 4.34.0 and the same pinned Mathlib revision in $66.321$ seconds.
All eight cover-stage dependency audits contain only \texttt{propext},
\texttt{Classical.choice}, and \texttt{Quot.sound}; the false arithmetic
control was rejected. Local and returned source hashes agree. The first
attempt failed on one ambiguous lemma name and remains recorded as a
failure. Both attempts used bounded memory and time and an independent
cloud deletion guard; cleanup evidence is recorded separately.

\textbf{Assessment.} Both optimization equalities and their common-cap
corollary are now formalized for the actual graphs. This closes the
universal multiplicity lemma, not merely another finite range of cases.
Uniform matching-class sizes and the single-neighborhood allocation
are also now verified. The independent two-type allocation and its
overlap estimate are now verified too, as is the coupled-color lower bound.
The full
two-type Tuza theorem and its numerical certificates are not yet fully
formalized; translating the full split-graph representation and handling
inactive vertices also remain obligations.

\paragraph{Sharpness.}
For packing, take $H=K_s$. Every triangle consumes at least one clique edge,
so $\tp\leq\binom{s}{2}$. Equality forces every packed triangle to be centered
at a new vertex: a clique triangle would consume three clique edges.
Such an equality packing is precisely a proper edge coloring of $K_s$ by
its centers. It exists at $m=h(s)$ and cannot exist at $m=h(s)-1$.
For covers, at $m=s-1$ the replacement argument gives $\tc=\binom{s}{2}$,
whereas at $m=s-2$ one may leave a star in the clique and remove the spoke
from its center to each new vertex, using $\binom{s}{2}-1$ edges.

\begin{corollary}
Every two-type instance has an equivalent instance for both $\tc$ and $\tp$
with at most $k+h(s)+h(t)\leq3k$ vertices. It suffices to truncate the two
types successively. Inactive types can be deleted altogether.
\end{corollary}
This does not bound $k$. It reduces multiplicities uniformly, rather than
verifying the conjecture for an additional finite clique size.

\section{An extremal lemma for two forbidden cliques}
For subsets $P,Q$ of a $k$-element set $C$, define
\[
 R(P,Q)=K_C-\big(E(K_P)\cup E(K_Q)\big).
\]
For a graph $R$, let $\operatorname{ex}_{\triangle}(R)$ be the maximum number
of edges in a spanning triangle-free subgraph of $R$.

\begin{lemma}\label{lem:twoholes}
For every $P,Q\subseteq C$,
\[
 \operatorname{ex}_{\triangle}(R(P,Q))=\mc(R(P,Q)).
\]
\end{lemma}
\begin{proof}
Partition $C$ into $A=P\setminus Q$, $B=Q\setminus P$, $D=P\cap Q$,
and $W=C\setminus(P\cup Q)$. In $R(P,Q)$, the sets $A,B,D$ are independent,
the pair $A,B$ is complete, $D$ has no neighbors in $A\cup B$, and $W$ is
a clique complete to all other vertices.

Take a triangle-free spanning subgraph $F$ with the maximum number of edges.
Within $A$, clone a vertex of maximum $F$-degree onto every other vertex of
$A$: give all of them its neighborhood. This is allowed by the host graph,
creates no triangle, and does not decrease the edge count. Perform the same
operation in $B$ and then $D$. Previously obtained twin classes remain twin
classes. Thus each of $A,B,D$ is homogeneous in $F$.

Next consider equivalence classes in $W$ under equality of open
$F$-neighborhoods. If two different classes have nonadjacent representatives
$v,w$, then $d_F(v)=d_F(w)$: otherwise cloning the lower-degree vertex onto
the higher-degree vertex would increase the edge count. Replace every
vertex in the class of $w$ by a clone of $v$. The two classes have no edges
between them, so the edge count is unchanged. No triangle is created and
no existing twin class is split, while these two classes merge. This also
preserves homogeneity of $A,B,D$. Repetition terminates with every two
distinct classes in $W$ adjacent. Since $F$ is triangle-free, there are
at most two such classes, say $W_0,W_1$.

The resulting graph is a blow-up of a triangle-free graph on at most five
vertices $A,B,D,W_0,W_1$ (empty classes are omitted). If there are five
nonempty classes, $W_0W_1$ is an edge, and $D$ can be adjacent only to
$W_0,W_1$. It cannot be adjacent to both, because that would make a triangle.
Thus $D$ has degree at most one in the quotient. The quotient has no odd
cycle: a triangle is excluded, and a five-cycle would have to pass through
$D$. It is therefore bipartite, as is its blow-up $F$.

The edge count of a bipartite subgraph is at most a maximum cut of the host.
Conversely every cut is a triangle-free subgraph. This proves equality.
\end{proof}

\section{An exact formula for the cover number}
Write $e_0=\binom{k}{2}$. The exact split-graph cover reduction of
Zeng~\cite{zeng}, or simply choosing the surviving neighborhood of each
independent-side vertex, gives
\begin{align}
 |E(G)|-\tc(G)
  &=\max_{\substack{L\subseteq K_C\\L\text{ triangle-free}}}
       \bigl(|E(L)|+p\alpha(L[S])+q\alpha(L[T])\bigr)\label{eq:oldcover}\\
  &=\max_{\substack{P\subseteq S\\Q\subseteq T}}
       \bigl(p|P|+q|Q|+\operatorname{ex}_{\triangle}(R(P,Q))\bigr).
       \label{eq:reorder}
\end{align}
Indeed, retained neighbors of any independent-side vertex must form an
independent set in $L$. For a fixed $L$, all copies of a type can choose the
same maximum independent set. Conversely, fixing $P,Q$ only forbids clique
edges internal to $P$ or $Q$; the remaining core must be triangle-free.
Lemma~\ref{lem:twoholes} replaces the last extremal problem by a maximum cut.

For nonnegative $\alpha,\beta,\gamma$ with sum at most $k$, set
$w=k-\alpha-\beta-\gamma$ and define
\begin{align}
 M_k(\alpha,\beta,\gamma)=
 \max_{\substack{\varepsilon_A,\varepsilon_B,\varepsilon_D\in\{0,1\}\\
                  z\in\{0,\ldots,w\}}}
 \bigl[&\alpha\beta\,\mathbf1_{\varepsilon_A\ne\varepsilon_B}
       +z(w-z)\nonumber\\
       &+d(w-z)+(\alpha+\beta+\gamma-d)z\bigr],\label{eq:M}\\
 d&=\varepsilon_A\alpha+\varepsilon_B\beta+\varepsilon_D\gamma.\nonumber
\end{align}
This is exactly $\mc(R(P,Q))$ when
$(|P\setminus Q|,|Q\setminus P|,|P\cap Q|)=(\alpha,\beta,\gamma)$.
To see this, a maximum cut may place every member of each false-twin class
$A,B,D$ on one side: with other vertices fixed, its cut contribution is
linear in the number placed on either side. Let $z$ vertices of $W$ be on
the same side as the $d$ counted vertices. The four terms in~\eqref{eq:M}
count $A$--$B$, $W$--$W$, and the two directions between $W$ and its complement.

For each of the eight assignments the part depending on $z$ is
\[
 wd+(k-2d)z-z^2.
\]
Only the nearest integers to $(k-2d)/2$, clipped to $[0,w]$, need be checked.
Thus no enumeration of core edge sets or cuts is required.

\begin{theorem}[Cover formula for every two-type instance]\label{thm:cover}
Let $\mathcal D$ consist of triples of nonnegative integers satisfying
\[
 \gamma+\max(0,\alpha-a)+\max(0,\beta-b)\leq u.
\]
Then
\[
 \boxed{\tc(G)=e_0+ps+qt-
  \max_{(\alpha,\beta,\gamma)\in\mathcal D}
   \bigl[p(\alpha+\gamma)+q(\beta+\gamma)
          +M_k(\alpha,\beta,\gamma)\bigr].}
\]
\end{theorem}
\begin{proof}
The displayed constraint is exactly the feasibility condition for
$P\subseteq S,Q\subseteq T$ with the three prescribed region sizes.
Their intersection uses $\gamma$ vertices of $S\cap T$. The two exclusive
parts must use at least $\max(0,\alpha-a)$ and $\max(0,\beta-b)$ further,
disjoint vertices of that intersection. This condition is also sufficient:
fill the exclusive parts from $S\setminus T,T\setminus S$ as far as possible,
and then use distinct remaining vertices of $S\cap T$.
In particular $\alpha+\beta+\gamma\leq a+b+u\leq k$.
Now apply~\eqref{eq:reorder}, Lemma~\ref{lem:twoholes}, and~\eqref{eq:M}.
\end{proof}
The formula has three bounded integer loops (each at most $k$), eight binary
assignments, and at most two choices for $z$ per assignment. This is an
$O(k^3)$ arithmetic-operation procedure; bit complexity also depends on the
input integers. It is an exact cover computation, not a packing formula.

\section{Two useful counterexamples to shortcuts}
\paragraph{Two optimal packings cannot simply be added.}
Let $C=\{u,v,a,b\}$, $S=\{u,v,a\}$, $T=\{u,v,b\}$, and $p=q=3$.
Each type considered with its own three-vertex clique admits three
edge-disjoint triangles. But $E(K_S)\cup E(K_T)$ has only five edges:
the two sets share $uv$. In the whole graph six packed triangles would
require all six core edges, with exactly one per triangle. The edge $ab$
cannot be used by an outside-centered triangle, so six are impossible.
Five are attained by
\[
 x_1uv,\quad x_2ua,\quad x_3va,\quad y_1ub,\quad y_2vb.
\]
The five core edges other than $ab$ cover all triangles. Therefore
$\tc(G)=\tp(G)=5$, whereas the two separate packing values sum to six.
This refutes the additive shortcut, not Tuza's conjecture.

\paragraph{The surviving whole graph need not be bipartite.}
Let $C=\{a,b,u\}$, $S=\{a,u\}$, $T=\{b,u\}$, and $p=q=2$.
Deleting $au,bu$ destroys all triangles, and the edge-disjoint triangles
$x_1au,y_1bu$ show $\tc=2$. There are eleven edges, so the maximum
triangle-free subgraph has nine edges. But the maximum cut has only eight:
if the three core vertices have the same color, all eight spokes cross;
otherwise two core edges cross, and the outside contribution is at most
six (or four when $u$ is alone on its side). Thus full-graph bipartization
costs three edges. Lemma~\ref{lem:twoholes} concerns the auxiliary core host,
and does not assert $\tc(G)=|E(G)|-\mc(G)$.

\section{An additive bound for an incomplete clique}
The following auxiliary inequality will absorb the part of the clique
left after assigning edges to outside-centered triangles. It does not
assume Tuza's conjecture for that remaining graph.

\begin{lemma}[Missing-edge allowance]\label{lem:allowance}
Let $H$ be any graph on $k$ vertices and let
$d=\binom{k}{2}-|E(H)|$. Then
\[
 \tc(H)\leq 2\tp(H)+d.
\]
\end{lemma}
\begin{proof}
Put $e=\binom{k}{2}$ and $v=\tp(K_k)$. The case $k\leq2$ is immediate.
Two elementary inequalities, valid for every graph, are
$\tc(H)\leq3\tp(H)$ and $\tc(H)\leq |E(H)|/2$.
For the first, the edges of a maximum triangle packing meet every triangle;
for the second, keep a cut containing at least half the edges. Consequently
\[
 \tc(H)-2\tp(H)\leq\tc(H)/3\leq(e-d)/6.
\]
If $d\geq e/7$, the last quantity is at most $d$, proving the claim.

Suppose now that $d<e/7$. Set $c=\lfloor k^2/4\rfloor$ and $\rho=c/e$.
A uniformly random balanced cut separates each pair with probability
$\rho$, so $\tc(H)\leq(1-\rho)(e-d)$. Deleting the $d$ missing edges
destroys at most $d$ triangles of a fixed packing in $K_k$, whence
$v\leq d+\tp(H)$. Thus it suffices to show
\begin{equation}\label{eq:small-defect}
 d+\tc(H)\leq e-c+\rho d\leq 2v.
\end{equation}
For $k\geq6$, the classical complete-graph packing formula, as stated in
\cite[Lemma 6]{bonamy}, implies
\[
 v\geq\frac{e-k/2-1}{3}.
\]
Indeed, the leave of an optimal packing has size $0,4,k/2$, or $k/2+1$
according to the residue of $k$ modulo six, and is at most $k/2+1$ in this
range. Moreover
\[
 \frac{2(e-k/2-1)}{3}-\left(e-\frac{6c}{7}\right)
 =\begin{cases}
 (2k^2-7k-28)/42,&k\text{ even},\\
 (2k^2-7k-37)/42,&k\text{ odd}.
 \end{cases}
\]
Both expressions are nonnegative: their smallest relevant arguments are
$6$ and $7$, and the quadratics increase thereafter. Since $d<e/7$,
$e-c+\rho d\leq e-6c/7\leq2v$, proving~\eqref{eq:small-defect}.

For $k=3,4$, the small-defect case forces $d=0$, and
$(e-c,v)=(1,1),(2,1)$. For $k=5$, it forces $d\leq1$; integrality gives
$d+\tc(H)\leq\lfloor4+3d/5\rfloor=4=2v$.
These small clique packings can be taken explicitly: one triangle for
$k=3,4$, and $123,145$ for $k=5$.
\end{proof}

\section{A complete allocation of the neighborhood edges}
For an arbitrary split graph, let $Q\subseteq E(K_C)$ be the set of core
edges belonging to at least one $K_{N(x)}$, $x\in I$. Following
\cite{zeng}, this is the neighborhood 2-shadow, here viewed as an edge set.
A \emph{full allocation} assigns every edge $uv\in Q$ to a common neighbor
$x\in I$, with edges assigned to each $x$ forming a matching. Equivalently,
it is an edge-disjoint packing of $|Q|$ outside-centered triangles.

\begin{theorem}[A full allocation suffices]\label{thm:full-allocation}
If a split graph admits a full allocation of $Q$, put $R=K_C-Q$ and
$D=|Q|$. Then
\[
 \tp(G)=D+\tp(R),\qquad \tc(G)=D+\tc(R),\qquad
 \tc(G)\leq2\tp(G).
\]
There is no restriction on the number of neighborhood types.
\end{theorem}
\begin{proof}
The allocated triangles can be combined with any packing in $R$, since
they use only edges of $Q$ and spokes. Conversely, in any packing each
triangle not contained in $R$ uses an edge of $Q$: this holds for core
triangles by definition of $R$, and for outside-centered triangles by
definition of $Q$. There are at most $D$ such triangles. This proves the
packing identity.

The set $Q$ together with a cover of $R$ covers every triangle of $G$.
Conversely, every cover must contain at least $\tc(R)$ edges of $R$ and
at least $D$ edges outside $R$ to hit the $D$ edge-disjoint allocated
triangles. This proves the cover identity. Finally
Lemma~\ref{lem:allowance}, applied to $R$, whose missing-edge count is $D$,
gives $\tc(R)\leq2\tp(R)+D$. Adding $D$ proves the conclusion.
\end{proof}

\begin{theorem}[A uniform multiplicity regime]\label{thm:saturation}
Consider any split graph. Suppose that each active neighborhood $S_i$
of size $s_i\geq2$ occurs at least $h(s_i)$ times. Then Tuza's inequality
holds. In particular, the two-type target holds whenever
\[
 p\geq h(s),\qquad q\geq h(t),
\]
for arbitrary $k,a,b,u,o$ and arbitrary overlap of the two neighborhoods.
\end{theorem}
\begin{proof}
Assign each edge of $Q$ to one type containing both its endpoints, choosing
only one type even when there are several possibilities. For each type,
restrict a fixed $h(s_i)$-edge-coloring of $K_{S_i}$ to its assigned edges.
Use a distinct outside vertex for each color. Within each center the edges
are a matching; between centers the spokes are distinct; and each core
edge was assigned only once. This is a full allocation, so
Theorem~\ref{thm:full-allocation} applies.
\end{proof}
This closes an infinite parameter regime rather than one more clique size.
The inequalities on multiplicities are sufficient, not necessary. They are
not asserted to be sharp for Tuza's inequality, even though the separate
compression thresholds in Theorem~\ref{thm:compression} are sharp.

\section{Partial allocations and their cost}
The preceding argument also measures what is missing when there are fewer
copies. Take any outside-centered packing, let $P\subseteq Q$ be its core
edges, and put $\ell=|P|$, $\delta=|Q\setminus P|$, $e=\binom{k}{2}$.

\begin{corollary}[Allocation deficit]\label{cor:deficit}
For every split graph and every such packing,
\begin{align*}
 \tc(G)-2\tp(G)&\leq\delta,\\
 \tc(G)-2\tp(G)&\leq\delta+\frac{e-7\ell}{6}.
\end{align*}
In particular, $7\ell\geq e+6\delta$ is a sufficient condition for Tuza's
inequality, even when the allocation is not full.
\end{corollary}
\begin{proof}
Let $H=K_C-P$. The given packing can be combined with any packing in $H$,
and $Q$ together with a cover of $H$ is a cover of $G$. Hence
\[
 \tc(G)-2\tp(G)
 \leq \delta-\ell+\tc(H)-2\tp(H).
\]
Lemma~\ref{lem:allowance} bounds the last difference by $\ell$, proving the
first claim. The elementary bound $\tc(H)-2\tp(H)\leq|E(H)|/6$ from the
same proof, with $|E(H)|=e-\ell$, proves the second.
\end{proof}

For two types the deficit has an explicit upper bound. Write
\[
 e_S=\binom{s}{2},\quad e_T=\binom{t}{2},\quad e_U=\binom{u}{2},\quad
 D=e_S+e_T-e_U,
\]
and $r_S=\min(p,h(s))/h(s)$, $r_T=\min(q,h(t))/h(t)$.
Choose uniformly a prescribed number of color classes from a fixed
complete-graph edge coloring for each type, independently between types.
Each edge of $K_S$ is selected with probability $r_S$, and each edge of
$K_T$ with probability $r_T$. If both types select an edge of $K_U$, retain
it only for the first type. Deleting such edges preserves every matching.
Linearity of expectation therefore gives a compatible packing of size at
least
\begin{equation}\label{eq:allocation-bound}
 L=\left\lceil r_Se_S+r_Te_T-r_Sr_Te_U\right\rceil.
\end{equation}
Thus $\delta\leq D-L$. If $L=D$, Theorem~\ref{thm:full-allocation}
applies. More generally, the condition $13L\geq e+6D$ implies Tuza's
inequality by Corollary~\ref{cor:deficit}. These are sufficient criteria;
failure of either criterion is not a counterexample to the conjecture.

\paragraph{A deterministic route to the independent-selection bound.}
The same bound has a short finite-sum proof useful for formalization.
Given $h$ nonnegative weights, choose $r$ smallest ones, with sum $W_A$;
let the sum of the unchosen weights be $W_B$. Comparing every chosen
weight to every unchosen weight gives
\[
 (h-r)W_A\leq rW_B,\qquad hW_A\leq r(W_A+W_B).
\]
This includes $r=0$ and $r=h$.

Put $p_0=\min(p,h(s))$ and $q_0=\min(q,h(t))$. Weight each $S$-color
by the number of its edges inside $U=S\cap T$, and choose the $p_0$
smallest weights. If $J_S$ is the selected edge set and
$d=|J_S\cap E(K_T)|$, the weights sum to $e_U$, so
$h(s)d\leq p_0e_U$. Next weight each $T$-color by the number of its
edges in $J_S$, and choose the $q_0$ smallest weights. If $J_T$ is its
selected edge set and $z=|J_S\cap J_T|$, these weights sum to $d$, so
\[
 h(t)z\leq q_0d,\qquad h(s)h(t)z\leq p_0q_0e_U.
\]
Keep all selected $S$-edges and only $J_T\setminus J_S$ for the second
type. Deleting edges preserves every matching. Uniform class sizes give
$|J_S|=p_0e_S/h(s)$ and $|J_T|=q_0e_T/h(t)$, and the compatible packing
has $|J_S|+|J_T|-z$ triangles. Its integer size is therefore at least
the ceiling in~\eqref{eq:allocation-bound}. The finite-sum argument,
palette choices and two-type graph construction are now formalized in
Lean in the integer form described below. This does not replace the
separate controlled-overlap argument for the coupled bound.

\paragraph{Lean verification of independent allocation.}
The module \path{lean-independent-allocation/IndependentAllocation.lean}
retains the preceding matching-class source verbatim, with one additional
Mathlib import for distributivity over finite sums. The theorem
\texttt{ordered\_subset} constructs an $r$-element subset by repeatedly
removing a least weight; it proves that every chosen weight is no larger
than any unchosen weight. The theorem \texttt{weighted\_subset} sums
these comparisons, including the boundary cases $r=0$ and $r=h$.
Two applications, in \texttt{select\_two\_palettes}, prove the product
bound on collisions without probability or an enumeration assumption.

The definition \texttt{twoAttached} is a Mathlib \texttt{SimpleGraph}
on an arbitrary old vertex type and two disjoint types of new centers.
Their neighborhoods are exactly $S$ and $T$; all new centers are mutually
nonadjacent. The theorem \texttt{merge\_palettes} assigns shared old
edges to the first type once and proves graph adjacency, edge-disjointness,
and the exact count $|P|=|J_S\cup J_T|$. It also exports the condition that
each produced triangle uses a new center, for later greedy completion.

For every pair of multiplicities, \texttt{independent\_allocation}
produces such a packing $P$ satisfying
\[
 h(t)p_0e_S+h(s)q_0e_T
 \leq h(s)h(t)|P|+p_0q_0e_U.
\]
Since $h(s),h(t)>0$, this is the denominator-cleared form of the
packing bound in~\eqref{eq:allocation-bound}. The Lean statement uses
the displayed integer inequality, rather than a rational ceiling.
For finite old vertex types, \texttt{independent\_allocation\_bound}
replaces $|P|$ by the actual attained packing optimum. The old graph
need only contain $S$ and $T$ as cliques; it need not be complete.
The result allows arbitrary overlap, identical or nested neighborhoods,
and zero multiplicities.

Azure job \texttt{31aa1dc2a94c4ffd} verified the complete source with Lean
4.34.0 and the pinned Mathlib revision in $79.186$ seconds. All seven
final audits contain only \texttt{propext}, \texttt{Classical.choice},
and \texttt{Quot.sound}; the false arithmetic control was rejected.
Local and returned source hashes agree. Two prior attempts failed on
elaboration and a missing import; both are retained as failures. A local
preparation superseded before launch created no cloud resources. Every
actual job had bounded memory and time and an independent cloud deletion
guard; cleanup evidence is recorded separately.

\textbf{Assessment.} One complete two-type graph-to-scalar packing lemma
is now kernel-checked, including the compatibility step and the center
condition needed later. The numerical formula and certificate domains
are unchanged, so existing numerical certificates were not rerun. The
coupled-palette lower bound below is now verified separately, using
coherent colors and controlled palette intersections. Core completion
is now verified too, with explicit preservation of the input packing.
The full two-type Tuza proof remains partly outside Lean.

\subsection{A shorter argument at half the multiplicity}
The structure of the residual core gives a stronger conclusion than the
general deficit estimate when $o\leq1$.

\begin{theorem}[Half the color threshold]\label{thm:half}
For a two-type split graph with $s,t\geq2$ and
$|C\setminus(S\cup T)|\leq1$, Tuza's inequality holds whenever
\[
 p\geq\lceil s/2\rceil,\qquad
 q\geq\lceil t/2\rceil.
\]
These thresholds equal $\lceil h(s)/2\rceil$ and $\lceil h(t)/2\rceil$.
There is no restriction on the sizes or overlap of the two neighborhoods.
More generally, in this setting it suffices that $2L\geq D$, with $L$
given by~\eqref{eq:allocation-bound}.
\end{theorem}
\begin{proof}
Put $R=K_C-Q$. Any outside-centered packing of size $L$ can be combined
with a packing in $R$, while $Q$ and a cover of $R$ cover all triangles.
Consequently
\[
 \tp(G)\geq L+\tp(R),\qquad
 \tc(G)\leq D+\tc(R).
\]
If $o=0$, then $R$ consists of the complete bipartite graph between
$A=S\setminus T$ and $B=T\setminus S$, with $U=S\cap T$ isolated.
Thus $\tc(R)=\tp(R)=0$.
If $o=1$, write $O=\{w\}$. The only triangles of $R$ are $wab$ with
$a\in A,b\in B$. A matching between $A,B$ gives $\min(a,b)$ packed
triangles; the edges from $w$ to the smaller side give a cover of the
same size. Hence again $\tc(R)\leq2\tp(R)$. This proves the conclusion
whenever $2L\geq D$.

Under the displayed multiplicity assumptions, $r_S,r_T\geq1/2$.
The expectation in~\eqref{eq:allocation-bound} is nondecreasing in each
probability, since its two partial derivatives are
$e_S-r_Te_U\geq0$ and $e_T-r_Se_U\geq0$. Therefore
\[
 L\geq \tfrac12e_S+\tfrac12e_T-\tfrac14e_U
      =\tfrac12D+\tfrac14e_U\geq\tfrac12D,
\]
as required.
\end{proof}
This proof uses only matching decompositions, averaging, and the displayed
description of the residual graph. It does not use
Lemma~\ref{lem:allowance} or the clique triangle-packing formula.
The multiplicity bound is sufficient and is not asserted to be sharp.

For example, take $a=b=u=10$, $o=0$, and $p=q=10$. Then $k=30$,
$s=t=20$, $h(s)=h(t)=19$, $D=335$, and
\[
 L=\left\lceil200-\frac{4500}{361}\right\rceil=188,
 \qquad \tc(G)\leq335<376\leq2\tp(G).
\]
These are certified bounds, not claimed exact optima. In fact the coarser
criterion $13L\geq e+6D$ narrowly fails here ($2444<2445$); the structural
argument above closes the case without improving the allocation bound.

\paragraph{Why the remaining allocation problem is real.}
Let $S=\{a,u,v,w\}$, $T=\{b,u,v,w\}$, and $p=q=2$.
Each type separately can pack four outside-centered triangles. Their
union cannot pack eight: equality would force all four centers to use
perfect matchings. The three shared vertices would then have total core
degree $12$, while $a,b$ contribute total degree $4$. Four distinct edges
inside $\{u,v,w\}$ would be required, but only three exist. Seven are
possible, for example
\[
 x_1au,\ x_1vw,\ x_2av,\ x_2uw,\ y_1bw,\ y_1uv,\ y_2bu.
\]
Formula~\eqref{eq:allocation-bound} gives
$\lceil4+4-4/3\rceil=7$, exactly the optimum for outside-centered
triangles in this example. This does not assert optimality among packings
that may also use core triangles.

\section{Using core triangles at low multiplicity}
The outside-allocation estimates ignore many available core triangles.
We now use the classical averaging method for clique packings, as in
Chahua and Guti\'errez~\cite[Corollary 8]{dense}, to obtain a complementary
regime. The stronger permutation-averaging argument remains as background;
the final proof uses the direct bound of Lemma~\ref{lem:sum-coloring}.
Both approaches allow triangles in the core and at outside centers.

Write $v_n=\tp(K_n)$, with $v_0=v_1=v_2=0$ and the classical values
from \cite[Lemma 6]{bonamy} for $n\geq3$.

\begin{lemma}[An elementary cyclic clique packing]\label{lem:cyclic-packing}
For every integer $n\geq3$, the complete graph $K_n$ has a triangle
packing of size at least $(n-1)(n-2)/6$.
\end{lemma}
\begin{proof}
Label its vertices by $\mathbb Z/n\mathbb Z$ and take every three-element
set $\{a,b,c\}$ with $a+b+c=0$. A pair $\{a,b\}$ determines its only
possible third vertex $c=-a-b$, so these triangles are edge-disjoint.
An uncovered edge must have $c=a$ or $c=b$. It therefore belongs to
\[
 \bigl\{\{v,-2v\}:v\in\mathbb Z/n\mathbb Z,\ v\ne-2v\bigr\}.
\]
There are at most $n-1$ such edges, because $v=0$ contributes no edge.
Thus the number of packed triangles is at least
\[
 \frac{\binom n2-(n-1)}3=\frac{(n-1)(n-2)}6.
\]
This construction requires no exact maximum-packing formula.
\end{proof}

\paragraph{Lean verification of the construction.}
The source \texttt{lean/CyclicPacking.lean} proves the construction first
for an arbitrary finite abelian group and then specializes to
$\mathbb Z/n\mathbb Z$. Its final theorem,
\texttt{TuzaCyclic.cyclic\_clique\_packing}, states that for every positive
$n$ there is a finite family $P$ of three-element vertex sets whose
two-element edge sets are pairwise disjoint and for which
\[
 \binom n2\leq3|P|+(n-1).
\]
This is the integral edge-count form of the bound above. The source
separately proves uniqueness of a triangle containing a prescribed edge,
the exact count $3|P|$ of covered edges, and the bound $n-1$ on uncovered
edges. It uses no enumeration over graph sizes or imported certificate
success assertion.

Azure job \texttt{5d63b3d4e13543f2} compiled the source with Lean 4.34.0
and pinned Mathlib dependencies. The final axiom report contains only
\texttt{propext}, \texttt{Classical.choice}, and \texttt{Quot.sound}; a
deliberately false arithmetic statement was rejected. The source contains
no admitted proof, custom axiom, unsafe declaration, or native decision
shortcut. The build used a 16 GiB whole-job memory limit, a 25-minute
runtime limit, and an independently armed cloud deletion guard. Logs,
source hashes, and resource-cleanup reports are retained under
\texttt{lean/azure/runs/}. This verifies the cyclic construction and
counting bound only. The stronger permutation-averaging lemma below and
the full two-type Tuza theorem are not yet Lean-verified. The final route
now uses the following separately verified direct construction.

\begin{lemma}[Triangle packing by sums of labels]\label{lem:sum-coloring}
Let $H$ be a finite simple graph on $n\geq1$ vertices, with $t(H)$
triangles. Its triangles can be assigned $n$ colors so that triangles
of the same color are edge-disjoint. Consequently,
\[
 \tp(H)\geq\left\lceil\frac{t(H)}n\right\rceil.
\]
In particular, applying this to $K_n$ gives the same guaranteed value
$g(n)=\lceil(n-1)(n-2)/6\rceil$ for $n\geq3$.
\end{lemma}
\begin{proof}
Label the vertices bijectively by $\mathbb Z/n\mathbb Z$. Color each
triangle by the sum of its three labels. If two triangles share an edge
with labels $a,b$ and have the same color, their third labels $c,d$
satisfy $a+b+c=a+b+d$, hence $c=d$. The triangles are therefore equal.
Each color class is an edge-disjoint packing. All $t(H)$ triangles lie
in these $n$ classes, so the largest class has at least $t(H)/n$
triangles; its size is an integer. For $K_n$, use
$t(K_n)=\binom n3$ and $\binom n3/n=(n-1)(n-2)/6$.
\end{proof}

\paragraph{Strict Lean statement and scope.}
The source \texttt{lean-sum-coloring/SumColoring.lean} proves the
statement for an arbitrary finite vertex type and a Mathlib
\texttt{SimpleGraph}. It returns a family $P$ of actual triangles,
proves pairwise disjointness of their two-element edge sets, and proves
$t(H)\leq |V(H)|\,|P|$. Its empty-vertex case is included. The proof first
treats any finite family of three-element vertex sets with injective
labels in a finite abelian group; it then uses the residue group of
order $|V(H)|$. A separate corollary applies this to all three-element
subsets and proves $\binom n3\leq n|P|$.

Azure job \texttt{86a6063f11ee4e5a} compiled the entire module with Lean
4.34.0. The three audited results, including
\texttt{TuzaSumColoring.graph\_triangle\_packing}, depend only on
\texttt{propext}, \texttt{Classical.choice}, and \texttt{Quot.sound}.
The deliberately false arithmetic control was rejected. No admitted
proof, custom axiom, unsafe declaration, or native decision shortcut is
used. This formal result replaces both the clique-packing input and the
weaker averaging input in the final argument; the other graph lemmas
and numerical certificates still require Lean formalization.

\begin{lemma}[Averaging a clique packing]\label{lem:averaging}
If a graph $H$ on $n\geq3$ vertices contains $t(H)$ triangles, then
\[
 \tp(H)\geq
 \left\lceil\frac{v_n t(H)}{\binom n3}\right\rceil
 \geq\left\lceil\frac{t(H)}n\right\rceil.
\]
\end{lemma}
\begin{proof}
Fix a packing of size $v_n$ in $K_n$ and relabel its vertices by a uniformly
random permutation. Keep only those triangles whose three edges belong to
$H$. The retained triangles are still edge-disjoint. Each triangle of the
fixed packing becomes a uniformly random three-element vertex set, so the
expected retained size is $v_n t(H)/\binom n3$. Some permutation attains at
least this expectation, and integrality gives the first ceiling.
Lemma~\ref{lem:cyclic-packing} gives
$v_n\geq(n-1)(n-2)/6=\binom n3/n$, proving the second inequality
without the classical maximum-packing formula. Equivalently, average
the explicit cyclic packing itself to obtain the weaker bound directly.
\end{proof}

For a split graph with $k$ core vertices and $m$ outside vertices of degrees
$d_1,\ldots,d_m$, put $n=k+m$ and $W=\sum_i d_i$. Its triangle count is
\[
 t(G)=\binom k3+\sum_{i=1}^m\binom{d_i}{2}.
\]
For two types this specializes to
$t(G)=\binom k3+p\binom s2+q\binom t2$.
In particular, the previously proved cover formula can be compared with
the following explicit lower bound, requiring no search for a packing:
\begin{equation}\label{eq:combined-bound}
 B=\max\left\{v_k,\ L,\
 \left\lceil\frac{v_{k+p+q}
  \left(\binom k3+p\binom s2+q\binom t2\right)}
  {\binom{k+p+q}3}\right\rceil\right\}\leq\tp(G).
\end{equation}
Here $L$ is the outside-allocation bound in~\eqref{eq:allocation-bound}.
Thus $\tc(G)\leq2B$ is a sufficient scalar criterion. The next section
exhibits a failure of its proposed universal version and strengthens the
packing construction.

\begin{theorem}[Few outside vertices]\label{thm:few-outside}
Let $G$ be any split graph whose clique part has size $k$ and whose active
independent part has size $m$. If
\[
 \boxed{k\geq4m+4,}
\]
then $\tc(G)\leq2\tp(G)$. The neighborhoods may be arbitrary; there is
no bound on their number of types and no lower bound on their sizes.
\end{theorem}
\begin{proof}
Delete inactive independent-side vertices, which belong to no triangle.
If $m=0$, use Tuza's inequality for complete graphs (also a consequence of
Lemma~\ref{lem:allowance} with no missing edges). Assume $m\geq1$, so $k\geq8$.

Partition the core into two parts whose sizes differ by at most one.
For each outside vertex, delete its edges to the part containing fewer of
its neighbors. The remaining graph is bipartite after also deleting all
edges internal to the two core parts. Therefore
\begin{equation}\label{eq:degree-cover}
 \tc(G)\leq\binom k2-\lfloor k^2/4\rfloor+\frac W2
 \leq\frac{(k-1)^2}{4}+\frac W2.
\end{equation}
The clique packing formula gives
\[
 2\tp(G)\geq2v_k\geq\frac{k^2-2k-3}{3}.
\]
Consequently~\eqref{eq:degree-cover} already proves the theorem when
\[
 W\leq R:=\frac{k^2-2k-15}{6}.
\]

Suppose $W>R$. By Cauchy--Schwarz and Lemma~\ref{lem:averaging},
\[
 2\tp(G)\geq\frac{1}{n}
 \left(\frac{k(k-1)(k-2)}3+\frac{W^2}{m}-W\right).
\]
Subtracting the last upper bound in~\eqref{eq:degree-cover}, and multiplying
by $n>0$, shows that it suffices to prove $F(W)\geq0$, where
\[
 F(W)=\frac{W^2}{m}-\frac{n+2}{2}W
       +\frac{k(k-1)(k-2)}3-\frac{n(k-1)^2}{4}.
\]
Since $m\leq(k-4)/4$ and $n\leq(5k-4)/4$, and $W\geq0$, we have
\[
 F(W)\geq f_k(W):=
 \frac{4W^2}{k-4}-\frac{5k+4}{8}W
 +\frac{k(k-1)(k-2)}3-\frac{(5k-4)(k-1)^2}{16}.
\]
This quadratic is strictly increasing on $[R,\infty)$, because
\[
 f_k'(R)=\frac{17k^2-16k-432}{24(k-4)}>0\qquad(k\geq8).
\]
For completeness, the remaining endpoint calculation has only nonnegative
coefficients after writing $h=k-8\geq0$:
\[
 36(k-4)f_k(R)
 =h^4+28h^3+241h^2+594h+252>0.
\]
Thus $F(W)\geq f_k(W)\geq f_k(R)>0$ in this second case, proving the result.
\end{proof}

\paragraph{A fixed witness for why mixing the triangles helps.}
Take $C=\{1,\ldots,11\}$, $S=\{1,\ldots,8\}$,
$T=\{4,\ldots,11\}$, and two copies of each type. The exact cover formula
gives $\tc(G)=35$, while the separate lower bounds are $v_{11}=17$ and
$L=16$. Neither is enough alone. However, there are
$\binom{11}{3}+4\binom82=277$ triangles and $v_{15}=35$, so averaging gives
\[
 \tp(G)\geq\left\lceil\frac{35\cdot277}{455}\right\rceil=22,
 \qquad 35\leq44\leq2\tp(G).
\]
The companion checker also constructs an explicit packing of size $23$:
start with the $35$ triples $\{i,j,i\mathbin{\mathrm{xor}}j\}$ on the nonzero
four-bit vectors, label the two $S$-centers $12,13$ and the two $T$-centers
$14,15$, and retain the triples that are triangles of $G$. It checks the
retained edges for disjointness and a cover of size $35$. This is one fixed
witness, not a verification of the general theorem.

\section{A failed comparison and a compatible mixed packing}
The scalar comparison in~\eqref{eq:combined-bound} is false in general.
A bounded Azure search found
\[
 (a,b,u,o,p,q)=(4,4,5,3,5,5).
\]
Here $k=16$, $s=t=9$, and $n=26$. The exact cover formula gives
$\tc(G)=75$, while all three entries in~\eqref{eq:combined-bound} equal
$37$:
\[
 v_{16}=37,\quad
 L=\left\lceil40-\frac{250}{81}\right\rceil=37,\quad
 \left\lceil\frac{104(560+360)}{2600}\right\rceil=37.
\]
This disproves that sufficient comparison; it does not disprove Tuza's
inequality. It tells us to construct compatible packings, rather than
merely take the largest of separate estimates.

\begin{lemma}[Keeping the residual core]\label{lem:wedges}
Take any outside-centered packing and let $P$ be its set of core edges.
Write $\ell=|P|$, $w(P)=\sum_{v\in C}\binom{d_P(v)}2$, and let $t(P)$
be the number of triangles in $P$. For $k\geq3$,
\begin{align*}
 \tp(G)&\geq\ell+
 \left\lceil\frac{v_k}{\binom k3}
  \left(\binom k3-(k-2)\ell+w(P)-t(P)\right)\right\rceil,\\
 \tp(G)&\geq\frac{(k-1)(k-2)}6+\frac{2\ell}{k}+\frac{2w(P)}{3k}.
\end{align*}
\end{lemma}
\begin{proof}
By inclusion--exclusion, the number of triangles in $K_C-P$ is
\[
 \binom k3-(k-2)\ell+w(P)-t(P).
\]
Each missing edge occurs in $k-2$ triples, each pair of adjacent missing
edges occurs in one triple, and triples having all three edges missing
must be subtracted once more. Apply Lemma~\ref{lem:averaging} to $K_C-P$
and combine its packing with the given outside-centered packing. They
share no edges, proving the first inequality.

For the second, use the weaker bound $\tp(K_C-P)\geq t(K_C-P)/k$.
Every triangle in $P$ contributes three distinct adjacent-edge pairs,
so $w(P)\geq3t(P)$. Substituting and simplifying gives the stated bound.
\end{proof}

The next construction makes the quantities in this lemma explicit. Its
palette is shared between the two types, so no core edge can be allocated
twice. Choose an integer
\[
 H\geq\max\{2,h(|S\cup T|),p+q\}.
\]
Properly edge-color $K_{S\cup T}$ using $H$ colors, allowing unused colors.
Assign distinct colors uniformly at random to the $p+q$ outside vertices.
At a vertex of type $S$ or $T$, retain just the edges of its assigned color
whose endpoints are both in its neighborhood. These edges are a matching,
so they give an outside-centered packing. Put
\begin{align*}
 \lambda_H&=\frac{p\binom s2+q\binom t2}{H},\\
 \omega_H&=\frac{3p(p-1)\binom s3+3q(q-1)\binom t3
       +pq\,u\big((s-1)(t-1)-(u-1)\big)}{H(H-1)}.
\end{align*}

\begin{theorem}[A shared palette and the residual core]\label{thm:joint}
For every two-type instance with $k\geq3$ and every such $H$,
\[
 \boxed{\tp(G)\geq J_H:=
 \left\lceil\frac{(k-1)(k-2)}6+\frac{2\lambda_H}{k}
                  +\frac{2\omega_H}{3k}\right\rceil.}
\]
There is no restriction on $o$, the overlap, or the relative multiplicities.
\end{theorem}
\begin{proof}
Each edge of $K_S$ receives an $S$-center with probability $p/H$, and
similarly for $T$. The two allocations are disjoint, even on shared edges.
Thus $\mathbb E\ell=\lambda_H$.

Two adjacent core edges have distinct colors. If both are assigned to
$S$-centers, the probability is $p(p-1)/(H(H-1))$; there are
$3\binom s3$ eligible unordered pairs. The corresponding term for $T$
is identical with $q,t$ substituted. For mixed assignments, choose the
shared endpoint in $S\cap T$, an $S$-neighbor and a $T$-neighbor, with the
neighbors distinct. There are
$u((s-1)(t-1)-(u-1))$ such choices, each with probability $pq/(H(H-1))$.
This counts the two possible type assignments separately when both are
eligible, as required. Hence $\mathbb E w(P)=\omega_H$.

Take expectations in the second inequality of Lemma~\ref{lem:wedges} and
use integrality of $\tp(G)$.
\end{proof}

\paragraph{An optional sharper version.}
For $H\geq3$, define
\[
 \xi_H=\frac{p(p-1)(p-2)\binom s3+q(q-1)(q-2)\binom t3
 +p(p-1)q\binom u2(s-2)+pq(q-1)\binom u2(t-2)}{H(H-1)(H-2)}.
\]
This equals $\mathbb E t(P)$. The first two terms assign all three triangle
edges to one type. For two $S$-assignments and one $T$-assignment, the
endpoints of the $T$-edge lie in $S\cap T$ and the third vertex lies in
$S$; this gives $\binom u2(s-2)$ choices. The other mixed term is analogous.
All three colors are distinct because the edge coloring is proper.
The first inequality of Lemma~\ref{lem:wedges}, with expectations taken
before the final ceiling, therefore yields the stronger bound
\begin{equation}\label{eq:exact-joint}
 \tp(G)\geq
 \left\lceil\lambda_H+\frac{v_k}{\binom k3}
  \left(\binom k3-(k-2)\lambda_H+\omega_H-\xi_H\right)\right\rceil.
\end{equation}

\section{Coupled colors and an elementary completion}
The moment calculation above is useful, but it is not needed to repair the
previous obstruction. We can coordinate the two local colorings more
efficiently, then complete the packing by a maximal packing in the core.

\begin{theorem}[Two local colorings with controlled overlap]\label{thm:coupled}
Let $s,t\geq2$, $d=\max(s,t)$, and $H=h(d)$. Put
\[
 \bar p=\min(p,H),\qquad \bar q=\min(q,H),\qquad
 z=\max(0,\bar p+\bar q-H).
\]
There is an outside-centered packing of size at least
\[
 \boxed{L_{\mathrm c}=\left\lceil
 \frac{\bar p\binom s2+\bar q\binom t2-z\binom u2}{H}
 \right\rceil.}
\]
If $s=t$ and $p+q\leq h(s)$, the maximum size of an outside-centered
packing is exactly
\[
 \boxed{(p+q)\lfloor s/2\rfloor.}
\]
These statements allow arbitrary overlap and arbitrary $o$.
The last identity concerns only outside-centered triangles, not the
maximum packing in the whole graph.
\end{theorem}
\begin{proof}
Fix an $H$-edge-coloring of $K_d$. Embed each of $S,T$ injectively in its
vertex set, using the same labels for vertices of $S\cap T$. This is
possible because the labels used by $S\setminus T$ and $T\setminus S$
may overlap: these are two separate embeddings, not an embedding of their
union. Pull back the coloring to each clique. Each local color class is
a matching, and a shared edge receives the same color in both cliques.
There is no requirement that the coloring be proper on $K_{S\cup T}$.

Identify the colors with $\mathbb Z/H\mathbb Z$. Fix palettes $A,B$
of sizes $\bar p,\bar q$ with $|A\cap B|=z$: choose $B$ from the
complement of $A$ until it is exhausted, then use the necessary colors
of $A$. For each $r\in\mathbb Z/H\mathbb Z$, use the palettes
\[
 A_r=\{c:c+r\in A\},\qquad B_r=\{c:c+r\in B\}.
\]
They retain the prescribed sizes and intersection. Assign their colors
to distinct outside centers within each type. Retain any edge selected
in both types only for the first type. All remaining triangles are
edge-disjoint: within a center their core edges form a matching,
between centers the spokes are distinct, and no core edge is retained
twice.

Let $J_r,K_r$ be the selected core-edge sets. As $r$ runs through all
$H$ shifts, each edge of $S$ is selected exactly $\bar p$ times and
each edge of $T$ exactly $\bar q$ times. Each common edge is selected
by both exactly $z$ times, because its color agrees in the two
neighborhoods. Summing the union identity gives
\[
 \sum_r |J_r\cup K_r|+z\binom u2
   =\bar p\binom s2+\bar q\binom t2.
\]
A shift maximizing the union has size at least its finite average.
Taking a ceiling proves the bound. This is a deterministic finite-sum
argument; no probability formalization is needed.

If $s=t$, every color class in the round-robin factorization of $K_s$ has
size $\lfloor s/2\rfloor$. When $p+q\leq H$, assign disjoint color sets to
the two types. No shared edge is selected twice, so every outside center
attains this many triangles. No outside center can support more: its
triangles must use disjoint pairs of its $s$ spokes. This proves exactness.
\end{proof}

\paragraph{Lean verification of the coupled lower bound.}
The module \path{lean-coupled-allocation/CoupledAllocation.lean} retains
\path{IndependentAllocation.lean} verbatim. The theorem
\texttt{common\_labels} first treats $|S|\leq|T|$: inject
$S\setminus T$ into $T\setminus S$, and fix every vertex of $T$.
The resulting map is injective separately on $S$ and $T$.
The theorem \texttt{coherent\_coloring} pulls back a complete-graph
coloring, treating the reverse inequality by symmetry. It produces
one edge-color function that is proper on each neighborhood, so
agreement on shared edges is automatic.

The theorem \texttt{minimal\_palettes} constructs the prescribed
palette sizes and minimal intersection. The theorem
\texttt{shift\_total} proves the finite counting identity for each
edge family; \texttt{select\_coupled} sums the union identity and takes
a maximizing shift. The earlier merging proof is generalized to arbitrary
color types in \texttt{merge\_general}. It constructs actual triangles
in the same \texttt{twoAttached} graph and checks adjacency,
edge-disjointness, cardinality, and that every triangle uses a new center.

For every pair of clique neighborhoods and every pair of multiplicities,
\texttt{coupled\_allocation} produces a packing $P$ with
\[
 \bar p\binom s2+\bar q\binom t2
 \leq H|P|+z\binom u2.
\]
This is the denominator-cleared integer form of the first bound in
Theorem~\ref{thm:coupled}; the formal statement uses natural-number
subtraction for $z=\max(0,\bar p+\bar q-H)$.
The construction includes zero multiplicities and equal or nested
neighborhoods, and allows an arbitrary ambient vertex type. The old
graph need only contain $S$ and $T$ as cliques. For finite old vertex
types, \texttt{coupled\_allocation\_bound} replaces $|P|$ by the actual
attained graph packing number. The separate upper bound proving the
equal-size optimum identity in Theorem~\ref{thm:coupled} remains a
written argument; that identity is not claimed as a formalized result here.

Azure job \texttt{659e03e8d7c248d8} verified the complete source in
$73.128$ seconds with Lean 4.34.0 and the pinned Mathlib revision.
All eight final dependency audits use only \texttt{propext},
\texttt{Classical.choice}, and \texttt{Quot.sound}. The false arithmetic
control was rejected, the successful log contains no proof errors or
\texttt{sorryAx}, and the returned source matches the local source.
The first attempt failed on a subtraction-injectivity lemma applied on
the wrong side and is retained as a failure. Both jobs had a 16 GiB
memory cap, a 25-minute worker limit and an independently armed cloud
deletion guard; cleanup reports are separate from proof success.

\textbf{Assessment.} Both allocation lower bounds used by the final route
are now formalized for actual graphs. This closes a universal dependency,
not merely an additional range of finite cases. For the earlier parameter
example $s=t=9$, $u=5$, $p=q=5$, the independent expression is
$2990/81$, with ceiling $37$, whereas the coupled expression is $350/9$,
with ceiling $39$. These are guaranteed sizes, not exact optima.
No numerical formula or certificate domain changed. The core-completion
bridge below is now verified too; it consumes the exported center condition
and preserves the input packing. The complete two-type Tuza proof and its
numerical certificates are not yet formalized.

For equal neighborhood sizes, the coupled bound is at least the old
independent-selection bound~\eqref{eq:allocation-bound}. Their positive
terms agree, while
\[
 \frac{z}{H}\leq\frac{\bar p\bar q}{H^2}.
\]
When $\bar p+\bar q>H$, this inequality is equivalent to
$(H-\bar p)(H-\bar q)\geq0$; otherwise it is immediate.
For unequal sizes no such dominance is asserted, so both bounds can be kept.

\begin{lemma}[Greedy completion in the core]\label{lem:completion}
Let $G$ be any split graph with a $k$-vertex clique and an outside-centered
packing of size $\ell$. Define
\[
 c_k=\binom k2-\lfloor k^2/4\rfloor.
\]
Then this particular packing can be extended to a packing of size at least
\[
 \boxed{\max\left\{\ell,
      \left\lceil\frac{c_k+2\ell}{3}\right\rceil\right\}.}
\]
There is no restriction on neighborhood types.
\end{lemma}
\begin{proof}
Delete the $\ell$ core edges used by the given packing. Repeatedly choose
a triangle in the remaining core and delete its three edges, until there
is no triangle. If $b$ triangles were chosen, the leftover graph has
$\binom k2-\ell-3b$ edges. Mantel's bound gives
\[
 \binom k2-\ell-3b\leq\lfloor k^2/4\rfloor,
 \qquad
 \ell+b\geq\left\lceil\frac{c_k+2\ell}{3}\right\rceil.
\]
The new triangles use neither selected core edges nor any spokes, so the
two packings combine. The lower bound $\ell$ is automatic.

For completeness, Mantel's bound here has a short proof. In a
triangle-free graph with $k>0$ vertices and $e>0$ edges, every edge $vw$
satisfies $d(v)+d(w)\leq k$. Summing over edges and applying
Cauchy--Schwarz gives
\[
 4e^2/k\leq\sum_v d(v)^2
   =\sum_{vw\in E}(d(v)+d(w))\leq ke.
\]
Hence $e\leq k^2/4$; the zero cases and the integer rounding are immediate.
\end{proof}

\paragraph{Lean verification of core completion.}
The module \path{lean-core-completion/CoreCompletion.lean} retains the
coupled-allocation source verbatim and adds the Mathlib Tur\'an import.
It first defines a graph from a finite family of two-element sets and
connects that representation to Mathlib's unordered graph edges.
The theorem \texttt{mantel\_family} applies
\texttt{SimpleGraph.CliqueFree.card\_edgeFinset\_le} at $r=2$, together
with \texttt{SimpleGraph.turanNumber\_two}, from the pinned library
source~\cite{mathlib-turan}. The degree-sum argument above remains an
elementary written proof; the Lean development reuses the library's
formal general Tur\'an proof. Mantel is neither an extra hypothesis nor
a custom axiom.

The theorem \texttt{residual\_triangle\_free} proves that any maximum
triangle packing leaves a triangle-free residual edge family. Otherwise
an additional edge-disjoint triangle could be inserted, contradicting
maximal cardinality. The used-edge count is exactly three times the
packing size, and \texttt{residual\_packing} combines this identity with
Mantel. These are abstract finite existence proofs; no enumeration is
run when proving the result for arbitrary vertex types.

The main theorem \texttt{core\_completion} assumes a graph on $V\oplus C$,
with finite complete core $V$ and independent outer vertices $C$. For
any finite packing $P$ whose triangles each contain an outer vertex,
it constructs a packing $R$ satisfying
\[
 P\subseteq R,\qquad
 2|P|+\binom{|V|}{2}\leq3|R|+\left\lfloor\frac{|V|^2}{4}\right\rfloor.
\]
Every input triangle uses exactly one core edge, and these edges are
distinct. The construction packs triangles into the unused core and
lifts them into the original graph. It checks adjacency, disjointness
between both families, exact union cardinality, and preservation of
\emph{the particular input packing}. It includes an empty core and
places no restriction on the number or sizes of outer neighborhoods.
The outer vertex type can even be infinite, provided the input packing
is finite.

For finite outer vertex types, \texttt{core\_completion\_bound}
compares the result to the actual attained graph optimum and also
exports $|P|\leq\tp(G)$. The corollary \texttt{two\_type\_completion}
applies directly to the \texttt{twoAttached} graph with a complete old
core, so both earlier allocation constructions supply its input. The
formal statements use the displayed integer inequality rather than a
rational ceiling. For $k=16$ and $|P|=39$, it yields $|R|\geq45$ while
retaining every initial triangle; this illustrates a universal result,
not an assertion about an exact optimum.

Azure job \texttt{054a1f953fcc47d0} verified the complete source in
$87.570$ seconds with Lean 4.34.0 and pinned Mathlib revision
\texttt{5ed296525643}. All seven final dependency audits use only
\texttt{propext}, \texttt{Classical.choice}, and \texttt{Quot.sound}.
The false arithmetic control was rejected; the successful log has no
proof errors or \texttt{sorryAx}, and local and returned source hashes
agree. The first attempt failed on two pair-membership lemma references
and an insertion-case elaboration; it remains a recorded failure.
Both jobs had a 16 GiB memory cap, a 25-minute worker limit and an
independently armed cloud deletion guard. Resource deletion is recorded
separately from proof success.

\textbf{Assessment.} This closes the graph-theoretic completion lemma
used by the main two-type route, including the required compatibility
with the previous allocation. It is not a claim that a new finite range
of instances has been checked. The scalar formulas and certificate
domains remain unchanged, so the numerical certificates were not rerun.
The explicit cut construction and its exact costs are now formalized
as described with Lemma~\ref{lem:explicit-cut}; rounding remains.
Graph representation, inactive vertices, counting/domain
reductions and sound certificate checking still remain; the complete
two-type Tuza result is not yet a single fully formalized theorem.

This completion needs neither the complete-graph triangle-packing formula
nor the moment calculations. With
\[
 L_* = \max(L,L_{\mathrm c}),\qquad
 M=\max\left\{L_*,\left\lceil\frac{c_k+2L_*}{3}\right\rceil\right\},
\]
we have $\tp(G)\geq M$. Thus $\tc(G)\leq2M$ is another sufficient scalar
criterion. It is not asserted to hold for every two-type graph.

\begin{corollary}[A total multiplicity condition]\label{cor:total-multiplicity}
Suppose $s=t\geq2$ and $o\leq1$. Put $m=p+q$, $H=h(s)$, and
$D=2\binom s2-\binom u2$. Tuza's inequality holds whenever
\[
 2\min(m,H)\lfloor s/2\rfloor\geq D.
\]
In particular it holds if $p+q\geq h(s)$, without a separate lower bound
on either $p$ or $q$.
\end{corollary}
\begin{proof}
Retain any $\min(m,H)$ outside centers and give them distinct colors as
in Theorem~\ref{thm:coupled}. This produces
$\min(m,H)\lfloor s/2\rfloor$ outside triangles. As in the proof of
Theorem~\ref{thm:half}, the residual graph $R=K_C-Q$ satisfies
$\tc(R)\leq2\tp(R)$ when $o\leq1$. Combine this packing with a packing
in $R$, and cover with $Q$ and a cover of $R$.
Finally $H\lfloor s/2\rfloor=\binom s2\geq D/2$.
\end{proof}

\begin{lemma}[Balancing multiplicities for a common lower bound]\label{lem:balance}
Fix $k,s=t,u$ and an integer $m\geq0$. Among graphs in this family with
$p+q=m$, the cover number is largest when
$(p,q)=(\lfloor m/2\rfloor,\lceil m/2\rceil)$.
\end{lemma}
\begin{proof}
For each triangle-free core $J$, the cover cost is
\[
 \binom k2-|E(J)|+p(s-\alpha(J[S]))+q(s-\alpha(J[T])).
\]
Its minimum over $J$ is the cover number. Extend the same finite minimum
of affine functions to real nonnegative $p,q$. On the line $p+q=m$ this
function is concave, and it is symmetric under $p\leftrightarrow q$,
because the two exclusive regions have equal sizes. Every integer between
$i$ and $m-i$ therefore has value at least the value at $i$. The two
central integers maximize it.
\end{proof}

After multiplicity compression, $p,q\leq H=h(s)$. In the equal-size
case, $L_{\mathrm c}$ depends on $p,q$ only through $m=p+q$; so do its
completion bound, $v_k$, and the whole-graph averaging bound. Consequently
their comparison with the exact cover formula needs only balanced
multiplicities. This is a reduction of that sufficient comparison, not
a claim that proving Tuza for the balanced graph alone would automatically
prove it for the others: the common packing lower bound must still suffice.

\subsection{A stronger and shorter proof for the scaled family}
\begin{lemma}[Lifting a split-graph certificate]\label{lem:lift}
Let a split graph $G$ have a clique of size $k$, a packing of $b$ triangles,
and a triangle cover of size $f$. For an integer $r\geq1$, replace every
core vertex by an $r$-vertex clique and every outside vertex by $r$
independent copies, replacing each old edge by the complete bipartite
graph between its two classes. Call the resulting split graph $G^{(r)}$.
Then
\[
 \tp(G^{(r)})\geq br^2,\qquad
 \tc(G^{(r)})\leq fr^2+k\binom r2.
\]
\end{lemma}
\begin{proof}
For each packed triangle with vertex classes $X,Y,Z$, index each class
by $\mathbb Z/r\mathbb Z$ and take
\[
 \{x_i,y_j,z_{i+j}\}\qquad(i,j\in\mathbb Z/r\mathbb Z).
\]
Every edge in each of the three complete bipartite blocks appears exactly
once: the missing index is determined by either addition or subtraction
modulo $r$. Different base triangles use disjoint edge blocks, so all
$br^2$ triangles are edge-disjoint.

For the cover, delete every edge block corresponding to a base-cover edge
and all edges internal to each of the $k$ core classes. A surviving
triangle would have to use three different classes and would project to
an uncovered base triangle. This is impossible.
\end{proof}

\begin{corollary}\label{cor:scaled-obstruction}
For every integer $r\geq1$, the two-type graph with
\[
 (a,b,u,o,p,q)=(4r,4r,5r,3r,5r,5r)
\]
satisfies Tuza's inequality. In fact
\[
 \tc(G)\leq83r^2-8r,\qquad \tp(G)\geq61r^2.
\]
\end{corollary}
\begin{proof}
Partition the core into a part consisting of all of $B=T\setminus S$,
all of $U=S\cap T$, and $r$ vertices of $A=S\setminus T$, and its
complement. These parts have sizes $10r$ and $6r$. Delete their internal
edges and, for every $S$-center, its $3r$ spokes to the smaller part.
All $T$-neighbors lie in the larger part. What remains is bipartite, and
the deletion count is
\[
 \binom{10r}{2}+\binom{6r}{2}+(5r)(3r)=83r^2-8r.
\]

The graph is precisely the lift of the $r=1$ graph from
Lemma~\ref{lem:lift}. The explicit packing below has $61$ triangles, so
its lift has $61r^2$ triangles. Alternatively its $75$-edge cover lifts
to the same bound $75r^2+16\binom r2=83r^2-8r$.
Finally $2(61r^2)-(83r^2-8r)=39r^2+8r>0$.
\end{proof}

At $r=1$, the coupled bound gives $L_{\mathrm c}=39$, and its greedy
completion bound is $45$. The moment bound~\eqref{eq:exact-joint} gives
$46$, while a new fixed construction supplies $61$ explicitly listed
edge-disjoint triangles and a $75$-edge cover. In this construction label
$U$ by $0,\ldots,4$ and the four exclusive vertices of each type by
$5,\ldots,8$ in its own local copy of $K_9$. Color an edge by the sum of
its labels modulo nine. Assign colors $0,1,2,3,8$ to the five $S$-centers
and colors $4,5,6,7,8$ to the five $T$-centers. The only common color is
eight, which appears on no edge of $K_U$. Thus all ten centers receive
four triangles, for a total of $40$. A fixed greedy pass through the core
adds $21$ more, leaving $17$ core edges and no triangle. The companion
checker verifies the full $61$-triangle packing and the cover directly;
$61$ is not claimed to be the maximum packing size.
To make the base packing explicit in these notes, use physical labels
$U=\{0,\ldots,4\}$, $A=\{5,\ldots,8\}$,
$B=\{9,\ldots,12\}$, and $O=\{13,14,15\}$; the local label of a
vertex $v\in B$ is $v-4$. The additional core triangles are
\[
\begin{array}{rrrr}
(0,5,13)&(0,6,7)&(0,14,15)&(1,5,12)\\
(1,6,13)&(2,5,11)&(2,12,13)&(3,10,13)\\
(3,11,14)&(3,12,15)&(4,9,10)&(4,11,13)\\
(4,12,14)&(5,8,9)&(5,10,14)&(6,8,10)\\
(6,9,11)&(7,8,11)&(7,9,13)&(7,10,15)\\
(8,13,14)&&&
\end{array}
\]
The cover deletes the edges within
$U\cup B\cup\{5\}$ and its core complement, and the three spokes to
$6,7,8$ from each $S$-center. It has $\binom{10}{2}+\binom62+15=75$
edges, and the remaining graph is bipartite.
The failure of the old bound at $r=1$ is not asserted to persist under every
scaling; the corollary proves the whole scaled family independently of it.

\section{Equal neighborhood sizes: a uniform large-core proof}
In this section $s=t\geq2$, $a=s-u$, $o=k-2s+u$, and $m=p+q$.
There is no restriction on the overlap, on $o$, or on the individual
multiplicities. Write $c=c_k$ and $v=v_k$. We use
\begin{equation}\label{eq:basic-estimates}
 \frac{k^2-2k}{4}\leq c\leq\frac{(k-1)^2}{4},\qquad
 2v\geq\frac{k^2-2k-3}{3},\qquad
 2\tp(G)\geq
 \frac{k(k-1)(k-2)/3+ms(s-1)}{k+m}.
\end{equation}
The last two bounds follow from the complete-graph packing formula and
Lemma~\ref{lem:averaging}; we will use them only for $k\geq35$.

\begin{lemma}[Two cover bounds]\label{lem:equal-covers}
For every equal-size instance,
\begin{align}
 \tc(G)&\leq c+m\max\{0,s/2-k/8,s-k/2\},\label{eq:equal-cut-cover}\\
 \tc(G)&\leq
 \begin{cases}
 c+(s-u)u,&s\leq k/2,\\
 \binom k2-s(k-s)+su-u^2,&s\geq k/2.
 \end{cases}\label{eq:equal-shadow-cover}
\end{align}
\end{lemma}
\begin{proof}
For the first bound, extend the cover function to real multiplicities as
in Lemma~\ref{lem:balance}. Concavity and symmetry give
$\tc(p,q)\leq\tc(m/2,m/2)$ for this weighted extension. We may therefore
construct covers with equal weights $m/2$.

Put $h=\lceil k/2\rceil$. Place all of $A=S\setminus T$ on one side of
a balanced core cut and all of $B=T\setminus S$ on the other. This is
possible because $a\leq\lfloor k/2\rfloor$. Placing the outside types
on opposite sides deletes $c$ core edges and has spoke cost $mu/2$.
Alternatively put both outside types opposite the larger core part.
Choose that part by taking vertices of $U$ first, then $A\cup B$, then
$O$. Its total number of neighbors, counted once for each type, is
$\min(2h,h+u,2s)$. The resulting spoke cost is
$\frac m2\max(0,2s-2h,2s-h-u)$. Combining these cuts, and using
\[
 \min\{u,\max(0,2s-2h,2s-h-u)\}
 \leq\max(0,2s-2h,s-h/2),
\]
gives $\tc(G)\leq c+m\max(0,s-h,s/2-h/4)$.
Since $h\geq k/2$, this proves~\eqref{eq:equal-cut-cover}.

For the second bound delete all edges of the neighborhood shadow $Q$
and keep a cut in $R=K_C-Q$, while retaining all spokes. Every outside
neighborhood is independent in the retained core, so this is a triangle
cover. Put $A,B$ on opposite sides of the cut, put $U$ on the same side
as $B$, and put $z$ vertices of $O$ on the side of $A$. The number of
retained edges is
\[
 a^2+ao+uz+z(o-z).
\]
If $s\leq k/2$, then $u\leq o$; choose an integer nearest to $(o+u)/2$.
The retained count is
$a^2+ao+\lfloor(o+u)^2/4\rfloor=\lfloor k^2/4\rfloor-au$.
If $s\geq k/2$, choose $z=o$, obtaining
$a^2+ao+uo=s(k-s)-su+u^2$.
Subtract these retained counts from $\binom k2$.
\end{proof}

\begin{theorem}[Equal neighborhood sizes above a fixed core threshold]\label{thm:equal-large}
If the two active neighborhoods have the same size and $k\geq35$, then
$\tc(G)\leq2\tp(G)$ for all multiplicities and overlaps.
\end{theorem}
\begin{proof}
First suppose $m\geq H=h(s)$. The distinct-color construction provides
an outside packing of size $\binom s2$, by retaining any $H$ centers.
Consequently the clique and completion bounds give
\[
 2\tp(G)\geq\max\left\{2v,
                     \frac{2c+2s(s-1)}3\right\}.
\]
If $s\leq k/2$, the second cover bound gives
$\tc(G)\leq c+s^2/4\leq c+k^2/16$. By~\eqref{eq:basic-estimates},
\[
 2v-(c+k^2/16)\geq\frac{k^2-8k-60}{48}>0.
\]
For $k/2\leq s\leq2k/3$, use $su-u^2\leq s^2/4$ in the cover bound.
The completion bound minus this cover bound is at least
\[
 D_1(s)=-\frac{k^2}{3}+ks-\frac{7s^2}{12}
                       +\frac k6-\frac{2s}{3}.
\]
This is concave in $s$, and its endpoint values are
\[
 D_1(k/2)=\frac{k(k-8)}{48},\qquad
 D_1(2k/3)=\frac{k(4k-15)}{54}.
\]
Both are positive. For $s\geq2k/3$, feasibility gives $u\geq2s-k\geq s/2$.
Since $su-u^2$ decreases on this range,
$\tc(G)\leq\binom k2-(k-s)^2$. The completion bound minus this bound is
at least
\[
 D_2(s)=\frac{2k^2}{3}-2ks+\frac{5s^2}{3}
                          +\frac k6-\frac{2s}{3}.
\]
It is increasing for $s\geq2k/3$ when $k\geq3$, and
$D_2(2k/3)=k(4k-15)/54>0$. This finishes the case $m\geq H$.

Now suppose $m<H$. Theorem~\ref{thm:coupled} gives an outside packing
of size $\ell=m\lfloor s/2\rfloor\geq m(s-1)/2$. Hence
\begin{equation}\label{eq:equal-low-bounds}
 2\tp(G)\geq2v,\qquad
 2\tp(G)\geq m(s-1),\qquad
 2\tp(G)\geq\frac{2c+2m(s-1)}3.
\end{equation}
When $s\leq k/4$, the first cover bound is $\tc(G)\leq c\leq2v$.

Suppose $k/4\leq s\leq3k/4$, so the cover bound is
$C=c+ms/2-mk/8$. If $m\geq(k+3)/3$, take one quarter of the first bound
in~\eqref{eq:equal-low-bounds} and three quarters of the third. This gives
\[
 2\tp(G)-C\geq\frac{v-c}{2}+\frac{m(k-4)}8
 \geq\frac{k-21}{24}>0.
\]
If $m\leq(k+3)/3$, use averaging instead. Define
\[
 F=\frac{k(k-1)(k-2)}3+ms(s-1)
 -(k+m)\left(\frac{(k-1)^2}{4}+\frac{ms}{2}-\frac{mk}{8}\right).
\]
It suffices that $F\geq0$. Completing the square yields
\begin{align*}
 48F&=P(k,m)+3m(4s-k-m-2)^2,\\
 P(k,m)&=4k^3-24k^2+20k-9mk^2+12mk-12m^2-24m-3m^3.
\end{align*}
For nonnegative $k,m$, $P$ is strictly decreasing in $m$, since
$\partial P/\partial m=-9k^2+12k-24-24m-9m^2<0$.
Thus it suffices to evaluate at $m=(k+3)/3$. Writing $x=k-35\geq0$ gives
\[
 9P\bigl(k,(k+3)/3\bigr)
 =8k^3-282k^2+117k-351
 =8x^3+558x^2+9777x+1294>0.
\]

Finally suppose $s\geq3k/4$, so the cover bound is
$C=c+m(s-k/2)$. If $m\geq(k+2)/2$, the outside packing alone gives
\[
 2\tp(G)-C\geq m(k/2-1)-c\geq(2k-5)/4>0.
\]
If $m\leq(k+2)/2$, put
\[
 F=\frac{k(k-1)(k-2)}3+ms(s-1)
 -(k+m)\left(\frac{(k-1)^2}{4}+m(s-k/2)\right).
\]
Now
\[
 12F=Q(k,m)+3m(2s-k-m-1)^2,\qquad
 Q(k,m)=k^3-6k^2+5k-3m^3-6m^2-6m.
\]
Again $Q$ decreases in $m\geq0$. At the endpoint, writing $y=k-17\geq0$,
\[
 8Q\bigl(k,(k+2)/2\bigr)
 =5k^3-78k^2-68k-120
 =5y^3+177y^2+1615y+747>0.
\]
Averaging proves the required inequality in the final case.
\end{proof}

\subsection{The genuinely finite remainder}
The preceding theorem leaves only $3\leq k\leq34$ for crossing active
types of equal size. Equal neighborhoods or a single active type are
threshold graphs, already covered by~\cite{bonamy}. For crossing types,
\[
 2\leq s\leq k-1,\quad
 \max(0,2s-k)\leq u\leq s-1,\quad
 2\leq m\leq2h(s)
\]
is a finite, exhaustive domain after multiplicity compression.
For each tuple the common packing lower bound is
\[
 K=\max\left\{v_k,L_{\mathrm c},
 \left\lceil\frac{c+2L_{\mathrm c}}3\right\rceil,
 \left\lceil\frac{v_{k+m}\bigl(\binom k3+m\binom s2\bigr)}{\binom{k+m}3}
 \right\rceil\right\}.
\]
It depends on the multiplicities only through $m$. By
Lemma~\ref{lem:balance}, it is enough to exhibit a cover of size at most
$2K$ for $p=\lfloor m/2\rfloor$, $q=\lceil m/2\rceil$.
The Azure-only worker \texttt{equal-finite/finite\_check.py} implements
exact integer arithmetic for this finite remainder and saves each cover
certificate. The bounded Azure run \texttt{eaf7da84d1a44ab7} completed all
$81\,576$ tuples in $3\,384$ profiles on 7 October 2026, with no failed
inequality. Every recorded cover is either a whole-graph cut cover or
the shadow together with the complement of a cut in $R$. The full exact
cover formula was never needed. Each certificate's region sizes and cost
were checked separately from its selection. All lower bounds use integer
arithmetic and ceilings. Six polynomial identities from the large-core
proof were also checked using exact rational coefficients.

The returned \texttt{case-certificates.jsonl} file has SHA-256
\begin{center}\small\ttfamily
0a76f1189f9dae0d5de5cc5f2d33f49bad17\\
1f9f88793a421b7b2c50ca50f962
\end{center}
(the two lines concatenate to one digest). The domain, source hashes,
execution result, and certificates are retained under
\texttt{equal-finite/azure/runs/eaf7da84d1a44ab7/}.
The run checked the finite inequalities, not the logical validity of all
written lemmas; no Lean verification or independent general-proof audit
is claimed. The large-core proof above does not depend on extrapolating
from this finite check.

\begin{theorem}[Two equal-size neighborhood types]\label{thm:equal-all}
Let $G$ be a split graph with any specified clique part $C$. Suppose the
active independent-side neighborhoods take at most two values, and that
these values have equal cardinality. Then
\[
 \boxed{\tc(G)\leq2\tp(G).}
\]
The clique size, the two multiplicities, the overlap, and the number of
core vertices outside the neighborhoods are unrestricted.
\end{theorem}
\begin{proof}
Delete inactive vertices. With at most one neighborhood type, use the
known threshold-graph theorem~\cite{bonamy}. Otherwise the two distinct
equal-size neighborhoods cross. For $k\geq35$, apply
Theorem~\ref{thm:equal-large}. For $k\leq34$, compress both multiplicities
to $H=h(s)$ by Theorem~\ref{thm:compression}. The resulting graph lies
in the finite domain above. Its packing number is at least the common
bound $K$, while Lemma~\ref{lem:balance} bounds its cover number by that
of the balanced instance with the same $k,s,u,m$. The saved certificate
for that instance gives a cover of size at most $2K$. Compression
preserves both parameters, so the original graph satisfies the inequality.
\end{proof}

The equal-size proof uses neither the exact cover formula of
Theorem~\ref{thm:cover} nor the moment calculations of
Theorem~\ref{thm:joint}. The cutoff $35$ is a sufficient proof threshold,
not a claimed sharp structural boundary or an upper limit on the theorem.

\section{A quantitative margin and transfer to unequal sizes}
The equal-size proof gives more than a nonnegative difference between
twice the packing number and the cover number. The following margin can
pay for changing spokes. Its constants are sufficient, not claimed optimal.

\begin{theorem}[A uniform equal-size margin]\label{thm:equal-margin}
Suppose the two neighborhood types have equal size $s\geq2$, the core
has $k\geq120$ vertices, and there are $m=p+q$ outside vertices. Then
\[
 \boxed{2\tp(G)-\tc(G)\geq
 \frac{k-21}{48}\min\left(k,\frac{2m}{3}\right).}
\]
There is no restriction on the multiplicities or on the overlap. The
statement also allows identical neighborhoods and a zero multiplicity.
\end{theorem}
\begin{proof}
Write $\Delta=2\tp(G)-\tc(G)$, $\rho=(k-21)/48$, and
$\mu=\min(k,2m/3)$. Retain $c,v,H=h(s)$ from the equal-size proof.
All cover bounds below are from Lemma~\ref{lem:equal-covers}.

If $m\geq H$, the first part of Theorem~\ref{thm:equal-large}'s proof
gives, uniformly in $s$,
\[
 \Delta\geq E(k):=\frac{k^2}{48}-\frac{k}{6}-\frac54.
\]
Indeed, this is the bound for $s\leq k/2$. In the next interval the
concave polynomial $D_1$ is bounded below by the smaller of
$k(k-8)/48$ and $k(4k-15)/54$, and in the last interval $D_2$ is at
least the latter value. Both values are at least $E(k)$ for $k\geq120$.
Since
\[
 E(k)-\rho k=\frac{13k-60}{48}>0,
\]
and $\mu\leq k$, the assertion follows in this case.

Now let $m<H$. In particular $m<k$, so $\mu=2m/3$.
The packing bounds~\eqref{eq:equal-low-bounds} are available.
If $s\leq k/4$, the cover bound is $c$, and
\[
 \Delta-\rho k\geq
 \frac{k^2-2k-15}{12}-\rho k
 =\frac{3k^2+13k-60}{48}>0.
\]

For $k/4\leq s\leq3k/4$, use the cover bound
$C=c+ms/2-mk/8$. If $m\geq3k/8$, the convex combination of the clique
and completion bounds used earlier gives
\begin{align*}
 \Delta-\rho\mu
 &\geq-\frac{k^2-2k+9}{24}+\frac{m(8k-15)}{72}\\
 &\geq\frac{k-72}{192}\geq0.
\end{align*}
If $m\leq3k/8$, use the averaging polynomial $P$ from the same proof.
It decreases with $m$, and
\[
 P(k,3k/8)=\frac{239k^3}{512}-\frac{339k^2}{16}+11k.
\]
This quantity is positive for $k\geq120$, as also follows from the
positive expression below. The completed square gives
$\Delta\geq P(k,3k/8)/(48(k+m))$. Since $k+m\leq11k/8$,
\[
 \Delta\geq A(k):=\frac{239k^2}{33792}-\frac{113k}{352}+\frac16.
\]
Here $\mu\leq k/4$. Writing $x=k-120\geq0$, we have
\[
 A(k)-\rho k/4
 =\frac{63k^2-7152k+5632}{33792}
 =\frac{63x^2+7968x+54592}{33792}>0.
\]
In particular $A(k)>0$ and $P(k,3k/8)=66kA(k)>0$, justifying the
denominator comparison.

For $s\geq3k/4$, use $C=c+m(s-k/2)$. If $m\geq5k/8$, the outside
packing alone gives
\begin{align*}
 \Delta-\rho\mu
 &\geq \frac{m(35k-51)}{72}-\frac{(k-1)^2}{4}\\
 &\geq\frac{31k^2+33k-144}{576}>0.
\end{align*}
If $m\leq5k/8$, the averaging polynomial $Q$ from the earlier proof
decreases with $m$, and
\[
 Q(k,5k/8)=\frac{137k^3}{512}-\frac{267k^2}{32}+\frac54 k.
\]
The completed square and $k+m\leq13k/8$ give
\[
 \Delta\geq B(k):=\frac{137k^2}{9984}-\frac{89k}{208}+\frac5{78}.
\]
Indeed, with $x=k-120$, the required positivity and margin follow from
\[
 B(k)-\rho(5k/12)
 =\frac{151k^2-7356k+1920}{29952}
 =\frac{151x^2+28884x+1293600}{29952}>0.
\]
Thus $B(k)>0$ and $Q(k,5k/8)=(39k/2)B(k)>0$, validating the denominator
comparison. Finally $\mu=2m/3\leq5k/12$, completing the proof.
\end{proof}

\begin{lemma}[Transferring a packing--cover margin across edge edits]
\label{lem:edit-transfer}
Let $G,H$ be graphs on the same vertex set. Put
$D=E(G)\setminus E(H)$ and $A=E(H)\setminus E(G)$. Then
\[
 \tc(G)\leq\tc(H)+|D|,\qquad
 \tp(G)\geq\tp(H)-|A|,
\]
and consequently
\[
 2\tp(G)-\tc(G)\geq2\tp(H)-\tc(H)-|D|-2|A|.
\]
\end{lemma}
\begin{proof}
Take a triangle cover of $H$, restrict it to $E(G)$, and add $D$.
A triangle of $G$ containing an edge of $D$ is covered by $D$; every
other triangle is a triangle of $H$ and meets the restricted cover.
For packing, remove from a maximum packing of $H$ all triangles that use
an edge of $A$. Since the packed triangles are edge-disjoint, at most
$|A|$ of them are removed. The surviving triangles form a packing of $G$.
Combining the two inequalities proves the last statement.
\end{proof}

\begin{theorem}[Unequal sizes with controlled editing cost]
\label{thm:unequal-transfer}
Let $2\leq s<t\leq k$, $k\geq120$, and put
\[
 \delta=t-s,\quad p_0=\min(p,h(s)),\quad q_0=\min(q,h(t)),
 \quad m=p_0+q_0,\quad w=\min(q_0,2p_0).
\]
Then, for every feasible overlap,
\begin{equation}\label{eq:unequal-margin}
 \boxed{2\tp(G)-\tc(G)\geq
 \frac{k-21}{48}\min(k,2m/3)-\delta w.}
\end{equation}
In particular, Tuza's inequality holds whenever
\begin{equation}\label{eq:weighted-transfer}
 \boxed{144\delta\min(q_0,2p_0)
 \leq(k-21)\min\bigl(3k,2(p_0+q_0)\bigr).}
\end{equation}
\end{theorem}
\begin{proof}
Compress the two multiplicities to $p_0,q_0$ using
Theorem~\ref{thm:compression}. Both $\tc$ and $\tp$ are preserved.
Write $G_0$ for the compressed graph. Since
$|T\setminus S|=t-|S\cap T|\geq t-s=\delta$, choose
$X\subseteq T\setminus S$ with $|X|=\delta$.

There are two ways to make the neighborhood sizes equal on the same
vertex set. Replacing $T$ by $T\setminus X$ deletes $\delta q_0$ spokes
from $G_0$ and adds no edges. The margin of the resulting equal-size
graph therefore transfers to $G_0$ with a loss of at most $\delta q_0$.
Alternatively replacing $S$ by $S\cup X$ adds $\delta p_0$ spokes and
deletes no edges. Its margin transfers with a loss of at most
$2\delta p_0$. Both equal-size graphs have the same $k$ and $m$, so
Theorem~\ref{thm:equal-margin} supplies the same lower margin for them.
Choose the smaller loss and apply Lemma~\ref{lem:edit-transfer}.
This proves~\eqref{eq:unequal-margin}; multiplication by $144$ gives
the sufficient condition~\eqref{eq:weighted-transfer}.
\end{proof}

\begin{corollary}[A size-only sufficient condition]\label{cor:near-equal}
For a split graph with at most two active neighborhood types, let
$\delta=|s-t|$. If
\[
 \boxed{k\geq\max(120,48\delta+21),}
\]
then $\tc(G)\leq2\tp(G)$, for arbitrary multiplicities and overlap.
\end{corollary}
\begin{proof}
Equal sizes follow directly from Theorem~\ref{thm:equal-margin}.
Otherwise relabel the types so that $s<t$. We have $q_0\leq h(t)\leq k$,
and
\[
 w=\min(q_0,2p_0)\leq\frac{2(p_0+q_0)}3.
\]
For the latter inequality, use $q_0\leq2p_0$ or $q_0\geq2p_0$
respectively. Thus $w\leq\min(k,2m/3)$. The assumed bound gives
$\delta\leq(k-21)/48$, so the right-hand side
of~\eqref{eq:unequal-margin} is nonnegative.
\end{proof}

For example, a size difference of two is covered for every $k\geq120$,
with no further constraint on the multiplicities. The weighted criterion
can also work far from equal sizes: at
$(k,s,t,u,p,q)=(120,30,90,10,1,89)$, its equal-size margin is
$123.75$ and its editing loss is $120$. This is only an illustration
of the sufficient criterion, not a claim that the example escapes all
previous sufficient criteria. At $(120,60,63,20,20,40)$ the margin is
$82.5$ and the loss is $120$, so this criterion fails; such failure does
not assert a counterexample to Tuza's conjecture.

The results in this section are analytic: they do not use the finite
certificate remainder for $k\leq34$. The bounded Azure job
\texttt{af130e010c9e4012} checked ten exact rational polynomial identities
used in the margin calculation and four prescribed scalar examples.
All identities passed and all four example outcomes matched their
specified expectations, including the failed sufficient criterion above.
The worker exited successfully and the returned source hashes matched.
Its reports are retained under
\texttt{stability-check/azure/runs/af130e010c9e4012/}.
These algebra checks are not an independent audit of the graph-theoretic
arguments and are not a Lean verification.

\section{A direct cover bound for unequal neighborhood sizes}

\begin{lemma}[An explicit cover from any core cut]\label{lem:explicit-cut}
Let $C$ be the $k$-vertex core, let $S,T\subseteq C$ be the two
neighborhoods, and let their multiplicities be $p,q\geq0$. For every
$A\subseteq C$, with $j=|A|$, there is a triangle edge cover $F$ with
\begin{align*}
 |F|={}&\binom j2+\binom{k-j}2\\
 &+p\min\{|S\cap A|,|S\setminus A|\}
  +q\min\{|T\cap A|,|T\setminus A|\}.
\end{align*}
Consequently the displayed expression is an upper bound on $\tc(G)$.
No optimality of this cut or of this particular cover is asserted.
\end{lemma}
\begin{proof}
Give the core parts $A$ and $C\setminus A$ different colors. Put each
outside type on the side containing fewer of its neighbors. Delete
all edges whose endpoints have the same color. Every triangle has a
same-color pair, so the deleted edges meet every triangle.

The deleted core edges are the pairs within $A$ and the pairs within
$C\setminus A$, counted by the two binomial terms. For an outside
vertex of type $S$, the deleted spokes are exactly its neighbors on
its assigned side, giving the smaller of the two intersection counts;
the argument for type $T$ is identical. Distinct outer vertices give
distinct spokes, and no spoke is a core edge. Thus the displayed sum
is the exact number of deleted edges.
\end{proof}

\paragraph{Lean verification of the cut construction.}
The module \path{lean-cut-covers/CutCovers.lean} retains the previous
core-completion source and appends the verified cover-compression
namespace before the new namespace \texttt{TuzaCutCovers}. It defines
\texttt{splitGraph} for a finite complete core $V$ and a finite outer
set $D$, with an arbitrary neighborhood $N(c)$ for each $c\in D$.
There is no restriction on the number of neighborhood types.
For a prescribed core part $A$ and a Boolean side $b(c)$ for each center,
\texttt{cutEdges} explicitly lists all deleted pairs and spokes.
The theorem \texttt{explicit\_cut} produces a cover of exact size
\[
 \binom{|A|}{2}+\binom{|V|-|A|}{2}
 +\sum_{c\in D}\begin{cases}
 |N(c)\cap A|,&b(c)=\mathrm{true},\\
 |N(c)\setminus A|,&b(c)=\mathrm{false}.
 \end{cases}
\]
The proof checks actual adjacency, the same-color pair in every
triangle, and every disjointness and injectivity fact needed for this
count. Empty parts, empty neighborhoods and zero multiplicities are
included. The explicit identity \texttt{splitGraph\_twoAttached}
connects the graph representations. The corollaries
\texttt{two\_type\_cut} and \texttt{two\_type\_cover\_bound} prove
Lemma~\ref{lem:explicit-cut} with an actual edge-family witness and the
previously verified attained optimum \texttt{coverNumber}.

Azure job \texttt{1b407cf947d14ecc} verified the complete cumulative
source in $78.454$ seconds with Lean 4.34.0 and pinned Mathlib revision
\texttt{5ed296525643}. All nine final dependency audits use only
\texttt{propext}, \texttt{Classical.choice} and \texttt{Quot.sound}.
The false arithmetic control was rejected; the successful log has no
proof errors or \texttt{sorryAx}, and returned and local sources agree.
Attempt \texttt{3bf77e7527254d80} failed at three representation
interfaces and remains a recorded failure. Both jobs used a 16 GiB cap,
a 25-minute worker limit and an independently armed cloud cleanup guard.
Deletion is checked separately from proof success.

\textbf{Assessment.} This supplies actual covers with exact costs for
arbitrary finite sizes, including the integer common cuts used in
Lemma~\ref{lem:common-cut}. It changes no scalar formula or certificate
domain. The full common-neighborhood rounding bound is now verified as described
with Lemma~\ref{lem:common-cut}. Balanced rounding is now verified too,
as described with Lemma~\ref{lem:separable-cover}, and the full separable
cover bound is now verified. Shadow-cover counts, the general split-graph
representation and inactive vertices, graph
identities/domain reductions and sound certificate evaluation remain.
The full two-type theorem is not yet completely formalized.

The transfer argument loses information by making the two sizes equal.
For covers, that loss can be avoided entirely. Define
\[
 f_k(d)=\max\{0,d/2-k/8,d-k/2\},\qquad
 c_k=\binom k2-\lfloor k^2/4\rfloor.
\]

\begin{lemma}[A separable cover bound for two arbitrary sizes]
\label{lem:separable-cover}
For two neighborhoods $S,T\subseteq C$ of any sizes $s,t$, with
multiplicities $p,q\geq0$,
\[
 \boxed{\tc(G)\leq c_k+p f_k(s)+q f_k(t).}
\]
No restriction on their intersection or on the unused core is needed.
\end{lemma}
\begin{proof}
We first allow a core vertex to be split fractionally between two parts,
each of total size $k/2$. Put
\[
 b_s=f_k(s),\quad b_t=f_k(t),\qquad
 r_s=s-b_s,\quad r_t=t-b_t,\quad u=|S\cap T|.
\]
The definition implies
\[
 0\leq b_s\leq s/2,\quad r_s\leq k/2,
 \qquad 0\leq b_t\leq t/2,\quad r_t\leq k/2,
 \qquad 2(b_s+b_t)\geq s+t-k/2.
\]
We construct a fractional cut and fix which part contains each outside
type, so that the expected spoke deletion costs per center are at most
$b_s,b_t$, respectively.

If $u\leq b_s+b_t$, retain $r_s$ units of $S$ on one side and $r_t$
units of $T$ on the opposite side. Disjoint choices of these retained
portions exist: first use exclusive vertices, then the common region.
The required common amounts are $(u-b_s)_+$ and $(u-b_t)_+$, whose sum
is at most $u$. Each retained portion has size at most $k/2$, so fill
the remaining spaces to obtain two equal parts. Put each outside type
opposite its designated retained portion. Its deleted spoke mass is
at most $b_s$ or $b_t$.

If $u\geq s+t-k/2-(b_s+b_t)$, place retained portions of sizes
$r_s,r_t$ on the same side. They can be chosen with union size
\[
 \max(r_s,r_t,r_s+r_t-u)\leq k/2,
\]
by overlapping them as much as their common region permits. Fill this
side to size $k/2$ and put both outside types on the opposite side.
Again the deleted spoke masses satisfy the desired bounds.
The two cases cover all values of $u$ because
$2(b_s+b_t)\geq s+t-k/2$.

To round the fractional cut, let $x_v\in[0,1]$ be the fraction of core
vertex $v$ assigned to the first side, so $\sum_v x_v=k/2$.
Keep the outside-type orientations fixed. The spoke cost is an affine
function $d+\sum_v w_v x_v$; the weights $w_v$ may be negative.
Let $A$ contain the $\lfloor k/2\rfloor$ smallest weights, and let
$b\notin A$ have the smallest remaining weight $a=w_b$.
For $k=0$ there is nothing to round. For $k>0$, the ordering gives
\[
 \sum_v w_v x_v-\sum_{v\in A}w_v
 \geq a\left(\sum_v x_v-|A|\right).
\]
Indeed, for $v\in A$ we have
$w_v(x_v-1)\geq a(x_v-1)$, whereas for $v\notin A$ we have
$w_vx_v\geq ax_v$.

For even $k$, the right-hand side is zero, so choose $A$.
For odd $k$, it is $a/2$: choose $A$ if $a\geq0$, and choose
$A\cup\{b\}$ if $a<0$. In the latter case the added cost is
$a\leq a/2$. In each case the chosen integral part has size
$\lfloor k/2\rfloor$ or $\lceil k/2\rceil$, and its spoke cost is
at most $p b_s+q b_t$.

Delete all within-part edges using the explicit cut cover of
Lemma~\ref{lem:explicit-cut}. Every such balanced cut has core cost
\[
 \binom{|A|}{2}+\binom{k-|A|}{2}
 =\binom k2-|A|(k-|A|)=c_k.
\]
Thus an actual triangle cover has size at most $c_k+p b_s+q b_t$,
as required.
\end{proof}

\paragraph{Lean checkpoint: balanced rounding with a cover witness.}
The rounding step above is now kernel-verified in
\texttt{lean-balanced-rounding/BalancedRounding.lean}.
The theorem \texttt{balanced\_round} applies to an arbitrary finite set,
arbitrary real weights (including negative ones), and every vector in
$[0,1]^k$ with total mass $k/2$. It returns a balanced subset whose
linear cost does not exceed the fractional cost. The theorem
\texttt{balanced\_core\_cost} proves the exact natural-number identity
for the constant $c_k$, including the empty core.

For fixed outside-type orientations, \texttt{spoke\_cost\_identity}
identifies the fractional spoke cost as an affine function.
The theorem \texttt{balanced\_cut\_witness} combines the rounding
with \texttt{explicit\_cut}: it returns an actual triangle edge cover,
its exact cardinality, and the bound by $c_k$ plus the fractional
spoke cost. The corollary \texttt{balanced\_cover\_bound} passes to
the attained optimum \texttt{coverNumber}.

Azure job \texttt{77db4a0436b04225} verified the complete cumulative
source on its first compiler attempt in $75.042$ seconds, excluding
provisioning and package installation. All nine final dependency audits
use only \texttt{propext}, \texttt{Classical.choice} and
\texttt{Quot.sound}. The false arithmetic control was rejected; the
successful log contains no proof errors or \texttt{sorryAx}, and the
returned source agrees with the local source. Lean 4.34.0 and Mathlib
revision \texttt{5ed296525643} remain pinned. The worker had a 16 GiB
cap, a 25-minute limit and an independently armed cloud deletion guard.
Resource deletion is checked separately from proof success.

\textbf{Assessment.} This closes the general rounding and graph-cover
construction, not merely finitely many instances. The fractional existence
step has subsequently been verified as described next. No scalar formula
or certificate domain changed, so the numerical certificates were not rerun.

\paragraph{Lean checkpoint: the full separable cover bound.}
The complete Lemma~\ref{lem:separable-cover} is now kernel-verified in
\texttt{lean-separable-cover/SeparableCover.lean}. The theorem
\texttt{fractional\_cut} proves existence under the abstract hypotheses
\[
 0\leq b_s\leq s,\quad 0\leq b_t\leq t,\quad
 s-b_s\leq k/2,\quad t-b_t\leq k/2,\quad
 s+t-k/2\leq 2(b_s+b_t).
\]
It constructs a function $x:C\to[0,1]$ of total mass $k/2$ and fixes
the outside-type orientations so that the deleted spoke masses are at
most $b_s,b_t$. Both overlap regimes are explicit. Uniform fractions
on $S\setminus T$, $T\setminus S$ and $S\cap T$ realize the prescribed
region masses. Their bounds and sums include empty regions. The lemma
\texttt{fill\_between} attains any total between two pointwise capacity
functions, by an explicit affine interpolation; its zero-denominator
case is handled separately.

The lemmas \texttt{penalty\_bounds} and \texttt{penalty\_gap} verify
these hypotheses for $b_s=f_k(s)$ and $b_t=f_k(t)$. The final theorem
\texttt{separable\_cover\_bound} invokes the already verified
\texttt{balanced\_cover\_bound}. Thus the graph-cover bound has no
unproved fractional-existence assumption.

The same source also verifies \texttt{four\_region\_cover}: any feasible
assignment of masses to all four regions totaling $k/2$ gives the
corresponding affine cover cost, with either choice of outer sides.
The theorem \texttt{greedy\_four} proves feasibility of greedily filling
four nonnegative capacities. The four specific piecewise formulas used
by the interval certificate still need to be connected to these lemmas.

Azure job \texttt{2e1913652b1c40c2} verified the cumulative source in
$76.560$ seconds, excluding provisioning and package installation.
All eleven final audits use only \texttt{propext},
\texttt{Classical.choice} and \texttt{Quot.sound}; the false arithmetic
control was rejected. Local and returned sources agree and the successful
log has no proof errors or \texttt{sorryAx}. Attempt
\texttt{d56a278a000c4b87} failed on definition elaboration, explicit casts
and a sum binder; its report is retained as a failure. Each worker had
a 16 GiB cap, a 25-minute limit and an independent cloud deletion guard.
Deletion is recorded separately from proof verification.

\textbf{Assessment.} The separable cover lemma is completely formalized
for arbitrary sizes, intersections and multiplicities. This closes a
universal lemma needed by both the analytic and finite branches. Shadow
covers, the particular fractional-cut formulas, graph bounds and
representations, scalar reductions, and sound verification of the two
certificates remain before the final theorem can be claimed in Lean.

For $s=t$, this recovers~\eqref{eq:equal-cut-cover} without exchanging
the two neighborhoods or balancing their multiplicities. The new lemma
removes the cover-symmetry obstacle completely. It does not assert that
different local matching capacities can all be attained simultaneously.

\begin{lemma}[A single type can pay for its separable cover cost]
\label{lem:single-type-certificate}
Let $k\geq35$, $2\leq d\leq k$, and $0\leq r\leq h(d)$.
Let $H$ consist of a $k$-vertex core and $r$ outside centers with one
common neighborhood of size $d$. Then
\[
 2\tp(H)\geq c_k+r f_k(d).
\]
\end{lemma}
\begin{proof}
The $r$ local colors give $r\lfloor d/2\rfloor$ outside triangles,
including at the endpoint $r=h(d)$. The clique, outside, completion,
and averaging bounds therefore give
\[
 2\tp(H)\geq\max\left\{2v_k,\ r(d-1),\
 \frac{2c_k+2r(d-1)}3,\
 \frac{k(k-1)(k-2)/3+r d(d-1)}{k+r}\right\}.
\]
Here is a self-contained scalar comparison, so this lemma does not
require the separate theorem for equal neighborhood sizes.
For $d\leq k/4$, the clique bound gives $2v_k\geq c_k$.
For $k/4\leq d\leq3k/4$ and $r\geq(k+3)/3$, take one quarter of
the clique bound and three quarters of the completion bound. Subtracting
$c_k+r(d/2-k/8)$ gives at least
\[
 \frac{v_k-c_k}{2}+\frac{r(k-4)}8
 \geq\frac{k-21}{24}>0.
\]
For $d\geq3k/4$ and $r\geq(k+2)/2$, the outside bound gives
\[
 r(d-1)-[c_k+r(d-k/2)]
 =r(k/2-1)-c_k\geq(2k-5)/4>0.
\]
It remains to justify the averaging comparison in the complementary
multiplicity ranges. For a nonnegative real $m$, define
\[
 F_d(m)=\frac{k(k-1)(k-2)}3+md(d-1)
 -(k+m)\left(\frac{(k-1)^2}{4}+m f_k(d)\right).
\]
The following identities hold in their indicated degree ranges:
\begin{align*}
 48F_d(m)&=P(k,m)+3m(4d-k-m-2)^2,
       &&k/4\leq d\leq3k/4,\\
 12F_d(m)&=Q(k,m)+3m(2d-k-m-1)^2,
       &&3k/4\leq d\leq k,\\
 P(k,m)&=4k^3-24k^2+20k-9mk^2+12mk-12m^2-24m-3m^3,\\
 Q(k,m)&=k^3-6k^2+5k-3m^3-6m^2-6m.
\end{align*}
Both $P,Q$ decrease with $m\geq0$. At their required endpoints,
\begin{align*}
 9P(k,(k+3)/3)&=8h^3+558h^2+9777h+1294>0,
       &&h=k-35\geq0,\\
 8Q(k,(k+2)/2)&=5j^3+177j^2+1615j+747>0,
       &&j=k-17\geq0.
\end{align*}
Thus $F_d(m)\geq0$ when $k/4\leq d\leq3k/4$ and
$m\leq(k+3)/3$, or when $d\geq3k/4$ and $m\leq(k+2)/2$.
Taking $m=r$ proves the remaining cases, since
$c_k\leq(k-1)^2/4$. The comparison includes $r=0$ and $r=h(d)$.
\end{proof}

\begin{theorem}[Two direct unequal-size regimes]\label{thm:direct-regimes}
Let $G$ have at most two active neighborhood types and a core of size
$k\geq35$. Compress their multiplicities to
\[
 p_0=\min(p,h(s)),\quad q_0=\min(q,h(t)),\quad m=p_0+q_0.
\]
Then Tuza's inequality holds if either
\[
 \boxed{\min(s,t)\leq k/4}
 \qquad\text{or}\qquad
 \boxed{3m\leq k+3.}
\]
Both statements allow arbitrary overlap; the first places no further
restriction on either multiplicity or on the larger neighborhood size.
\end{theorem}
\begin{proof}
Compression preserves both parameters, so work with $p_0,q_0$.
Relabel so that $s\leq t$.

If $s\leq k/4$, then $f_k(s)=0$, and
Lemma~\ref{lem:separable-cover} gives
$\tc(G)\leq c_k+q_0 f_k(t)$, independently of $p_0$.
The subgraph on the core and the $q_0$ centers of type $T$ is the graph
of Lemma~\ref{lem:single-type-certificate}. Its packing is also a
packing of $G$, so $2\tp(G)\geq c_k+q_0 f_k(t)\geq\tc(G)$.

Now suppose $3m\leq k+3$. The previous case permits us to assume
$s,t>k/4$; $m=0$ is just a clique. Put
\[
 F_d=\frac{k(k-1)(k-2)}3+m d(d-1)
 -(k+m)\left(\frac{(k-1)^2}{4}+m f_k(d)\right).
\]
For $k/4<d\leq3k/4$, the completed-square identity in the proof of
Lemma~\ref{lem:single-type-certificate} gives $48F_d=P(k,m)+3m(4d-k-m-2)^2\geq0$, since
$m\leq(k+3)/3$ and $k\geq35$.
For $d\geq3k/4$, the corresponding identity is
$12F_d=Q(k,m)+3m(2d-k-m-1)^2\geq0$.
Here we use $m\leq(k+3)/3\leq(k+2)/2$ and the positivity bound
for $Q$ proved in the same lemma (already valid for $k\geq17$).

Average $F_s,F_t$ with weights $p_0/m,q_0/m$. Their weighted sum is
\begin{align*}
 &\frac{k(k-1)(k-2)}3+p_0s(s-1)+q_0t(t-1)\\
 &\hspace{8mm}-(k+m)\left(\frac{(k-1)^2}{4}
                  +p_0 f_k(s)+q_0 f_k(t)\right)\geq0.
\end{align*}
The first line divided by $k+m$ is a lower bound for $2\tp(G)$
by clique averaging. The last parenthesis is an upper bound for
$\tc(G)$ by Lemma~\ref{lem:separable-cover} and
$c_k\leq(k-1)^2/4$. This proves the second assertion.
\end{proof}

The first new regime can have arbitrarily large relative size imbalance;
for example it covers all feasible pairs with $s\leq k/4$ and
$t$ close to $k$, regardless of the multiplicities. The second replaces
the earlier sufficient bound $m\leq(k-4)/4$ by
$m\leq(k+3)/3$ for two types and $k\geq35$.
Both are analytic deductions and use no finite remainder certificates.

\subsection{A shadow cover for unequal sizes}
\begin{lemma}[An explicit cut outside the neighborhood shadow]
\label{lem:unequal-shadow}
Suppose $s\leq t$ and put $u=|S\cap T|$. For arbitrary multiplicities,
\[
 \tc(G)\leq
 \begin{cases}
 c_k+(s-u)u,&t\leq k/2,\\
 \binom k2-t(k-t)+su-u^2,&t\geq k/2.
 \end{cases}
\]
The two expressions agree when $t=k/2$.
\end{lemma}
\begin{proof}
Write $A=S\setminus T$, $B=T\setminus S$, $U=S\cap T$ and
$O=C\setminus(S\cup T)$, with sizes $a=s-u$, $b=t-u$, $u$, and
$o=k-s-t+u$. Delete the entire neighborhood shadow
$E(K_S)\cup E(K_T)$. This destroys all outside-centered triangles.
In the residual core, keep a cut with $A$ on one side, $B\cup U$
on the other, and $z$ vertices of $O$ on the $A$ side. Its edge count is
\[
 ab+a(o-z)+tz+z(o-z)
 =a(k-s)+(k-2a)z-z^2.
\]
Since $a\leq b$, the real maximizer $z_*=k/2-a$ is nonnegative.
Also $a+o=k-t$, so $z_*\leq o$ exactly when $t\leq k/2$.
In that case choose an integer $z$ nearest to $z_*$; the retained
edge count is $\lfloor k^2/4\rfloor-au$. If $t\geq k/2$, choose
$z=o$ instead, retaining $ab+to=t(k-t)-su+u^2$ edges.
The surviving graph has a bipartite core and no outside-centered
triangle. The deleted core edges therefore form a triangle cover.
Subtract the retained edge count from $\binom k2$.
These explicit cuts are not asserted to be maximum cuts.
\end{proof}

\begin{corollary}[Two neighborhoods no larger than half the core]
\label{cor:two-half}
If $\max(s,t)\leq k/2$, then $\tc(G)\leq2\tp(G)$,
with no restriction on multiplicities or overlap.
\end{corollary}
\begin{proof}
The preceding lemma gives
\[
 \tc(G)\leq c_k+(s-u)u\leq c_k+s^2/4\leq c_k+k^2/16.
\]
The classical clique-packing bound and $c_k\leq(k-1)^2/4$ imply
\[
 2\tp(G)-\tc(G)\geq2v_k-c_k-k^2/16
 \geq\frac{k^2-8k-60}{48}>0 \qquad(k\geq13).
\]
Thus a packing using only core triangles already suffices.
For $6\leq k\leq12$, retain the integer bounds. The same argument
works by the following values from the complete-graph packing formula:
\[
\begin{array}{c|rrrrrrr}
 k&6&7&8&9&10&11&12\\\hline
 c_k+\lfloor\lfloor k/2\rfloor^2/4\rfloor&8&11&16&20&26&31&39\\
 2v_k&8&14&16&24&26&34&40
\end{array}
\]
For $k\leq3$ there is no active outside neighborhood, so the graph's
triangles lie in its clique. For $k=4,5$, an active neighborhood has
size two. With at most one distinct active neighborhood, the shadow
bound is $c_k\leq2v_k$. With two distinct active neighborhoods, choose
one center from each and pack their two triangles: their core edges
are different. When $k=4$ this suffices, since the shadow bound is at
most $c_4+1=3$. When $k=5$, at least eight core edges remain after these
two triangles, more than Mantel's bound of six. A third, core triangle
can be added, while the shadow cover has size at most $c_5+1=5$.
These small cases are direct arguments, not a computational remainder.
\end{proof}

\begin{corollary}[A large exclusive part]\label{cor:large-exclusive}
If $k\geq35$ and $|T\setminus S|\geq k/2$, then Tuza's inequality
holds for arbitrary multiplicities. In particular it holds whenever
$|s-t|\geq k/2$.
\end{corollary}
\begin{proof}
After compression, make a fractional balanced cut whose first side
of size $k/2$ lies entirely in $T\setminus S$. Put the $S$ centers on
that first side and the $T$ centers on the other side. No $S$ spoke is
deleted, and the deleted $T$ spoke mass is $q(t-k/2)$. Balanced rounding
therefore gives
\[
 \tc(G)\leq c_k+q(t-k/2)\leq c_k+q f_k(t).
\]
Lemma~\ref{lem:single-type-certificate} supplies a packing in the induced
one-type graph that pays for this bound. Finally,
$t-s\geq k/2$ implies $|T\setminus S|\geq t-s\geq k/2$; interchange
the types if needed.
\end{proof}

\begin{lemma}[An unbalanced cut inside the common neighborhood]
\label{lem:common-cut}
Put $m=p+q$, $u=|S\cap T|$, and $E=\binom k2$. For every
$0\leq\ell\leq u$,
\[
 \tc(G)\leq E+ps+qt-\ell(k+m)+\ell^2+\tfrac14.
\]
In particular choose $\ell=\min(u,(k+m)/2)$.
\end{lemma}
\begin{proof}
Put $B=k+m$, $Q(x)=x^2-Bx$, and choose
$j=\min(u,\lfloor B/2\rfloor)$. Make one core part a $j$-element
subset of $S\cap T$ and place all outside centers on the other side.
The explicit cut construction gives a triangle cover $F$ with
\[
 |F|=\binom j2+\binom{k-j}2+p(s-j)+q(t-j)
     =E+ps+qt+Q(j).
\]
If $u\leq\lfloor B/2\rfloor$, then $j=u$ and, for every $\ell\leq u$,
\[
 Q(\ell)-Q(j)=(u-\ell)(B-u-\ell)\geq0.
\]
Otherwise $j=\lfloor B/2\rfloor$. Since $B$ is an integer,
$0\leq B-2j\leq1$, and completing the square gives
\[
 Q(j)=-\frac{B^2}{4}+\frac{(B-2j)^2}{4}
 \leq-\frac{B^2}{4}+\frac14
 \leq Q(\ell)+\frac14.
\]
Thus the same cover $F$ satisfies the asserted bound for all
$0\leq\ell\leq u$. The minimum of $Q$ on this interval occurs at
$\ell=\min(u,B/2)$, yielding in particular
\[
 \tc(G)\leq E+ps+qt-\frac{B^2}{4}
           +\max(0,B/2-u)^2+\frac14.
\]
\end{proof}

\paragraph{Lean verification of the full common-neighborhood bound.}
The module \path{lean-common-rounding/CommonRounding.lean} retains
\path{lean-cut-covers/CutCovers.lean} verbatim, adds real-number and
arithmetic tactic imports, and appends namespace
\texttt{TuzaCommonRounding}. The theorem \texttt{commonCost\_real}
proves the identity between the natural-number edge count and its real
quadratic expression. The assumptions $j\leq k,s,t$ explicitly justify
every conversion of truncated natural subtraction to real subtraction.
The formula for $\binom n2$ is proved by induction from the binomial
recurrence, including $n=0$ and $n=1$.

The theorem \texttt{rounded\_quadratic} proves the two-case argument
above with $j=\min(u,\lfloor B/2\rfloor)$. The theorem
\texttt{rounded\_integer\_min} also proves $Q(j)\leq Q(r)$ for every
integer $0\leq r\leq u$. This minimizes the cost within the specified
common-neighborhood cut family; it does not assert that the resulting
cover is a minimum cover of the graph.

The theorem \texttt{common\_cut\_witness} returns an actual cover $F$,
its exact integer cost and a bound valid \emph{simultaneously for every
real $\ell\leq u$}. In particular it covers the entire interval in
Lemma~\ref{lem:common-cut}. The theorem \texttt{common\_cover\_bound}
passes to the attained graph optimum \texttt{coverNumber}, while
\texttt{common\_cover\_optimized} gives the clipped-square expression
at the end of the proof. Empty cores, empty neighborhoods and zero
multiplicities are included. No averaging theorem or numerical checker
output is used as a proof assumption.

Azure job \texttt{6f55a192e642448b} verified the cumulative source on its
first actual Lean compiler attempt. The build and audit took $67.403$
seconds, excluding provisioning and package installation. All eight
final dependency audits use only \texttt{propext},
\texttt{Classical.choice} and \texttt{Quot.sound}. The false arithmetic
control was rejected; the successful log contains no proof errors or
\texttt{sorryAx}, and local and returned source hashes agree. Lean 4.34.0
and Mathlib revision \texttt{5ed296525643} are pinned.

Earlier job \texttt{69df8614cfa84164} failed during Ubuntu package
installation, before Lean ran. Its environment-failure report is retained;
the bootstrap now has bounded package-download retries and timeouts.
Both jobs had a 16 GiB worker cap, a 25-minute limit and an independently
armed cloud deletion guard. Resource deletion is checked separately from
proof success.

\textbf{Assessment.} The whole common-neighborhood cover lemma, including
the $1/4$ allowance and its optimized form, is now kernel-checked. The
written proof has been simplified to the same deterministic argument.
No scalar formula or certificate domain changed, so no numerical
certificate replay was needed. Balanced rounding is now verified as well.
The full separable cover bound is now verified too. Shadow covers,
the specific fractional-cut formulas, the general split-graph representation and
inactive vertices, counting/domain reductions and sound certificate
evaluation still remain. The full two-type Tuza theorem is not yet
completely formalized.

\begin{corollary}[A common-neighborhood criterion]
\label{cor:common-cut}
Suppose $k\geq34$ and $p+q=m\leq2k$. Then Tuza's inequality holds if
\[
 \boxed{4(k+m)(k+m-2u)_+^2\leq k(k^2+3m^2).}
\]
After multiplicity compression, the condition $m\leq2k$ is automatic.
\end{corollary}
\begin{proof}
Use $x=s/k$, $y=t/k$, $\alpha=p/k$, $\beta=q/k$, $\mu=m/k$,
$\rho=1/k$, and put $e=((1+\mu)/2-u/k)_+$.
The preceding cover bound, divided by $k^2$, is
\[
 C=\frac{(1-\rho)^2}{4}+\alpha x+\beta y
                      -\frac\mu2-\frac{\mu^2}{4}+e^2.
\]
Let $A$ be the normalized averaging lower bound on $2\tp(G)$.
Completing squares gives the identity
\begin{align*}
 (1+\mu)(A-C)
 &=B(\mu,\rho)-(1+\mu)e^2\\
 &\quad+\alpha\left(x-\frac{1+\mu+\rho}{2}\right)^2
       +\beta\left(y-\frac{1+\mu+\rho}{2}\right)^2,\\
 B(\mu,\rho)
 &=\frac1{12}+\frac{\mu^2}{4}
       -\frac{(1+\mu^2)\rho}{2}
       +\left(\frac5{12}-\frac\mu2\right)\rho^2.
\end{align*}
For $0\leq\mu\leq2$ and $0\leq\rho\leq1/34$,
\[
 B(\mu,\rho)\geq
 \frac1{12}-\frac1{68}-\frac7{12\cdot34^2}
             +\left(\frac14-\frac1{68}\right)\mu^2
 =\frac{315}{4624}+\frac4{17}\mu^2
 \geq\frac{1+3\mu^2}{16}.
\]
The stated criterion is exactly
$16(1+\mu)e^2\leq1+3\mu^2$. Thus $A\geq C$.
\end{proof}

\subsection{The remaining explicit comparison}
For a compressed two-type instance, put $E=\binom k2$,
$R=K_C-(E(K_S)\cup E(K_T))$, and retain the independent and coupled
outside bounds $L,L_{\mathrm c}$. Define
\begin{align*}
 C_{\mathrm{dir}}&=\min\left\{
 \left\lfloor c_k+p f_k(s)+q f_k(t)\right\rfloor,\ E-\mc(R)\right\},\\
 B(r,z)&=\left\lceil\frac{v_{k+r+z}
       \bigl(\binom k3+r\binom s2+z\binom t2\bigr)}{\binom{k+r+z}3}\right\rceil,\\
 L_*&=\max(L,L_{\mathrm c}),\\
 K_{\mathrm{dir}}&=\max\left\{v_k,L_*,
  \left\lceil\frac{c_k+2L_*}{3}\right\rceil,
  B(p,q),B(p,0),B(0,q)\right\}.
\end{align*}
The two one-type averaging bounds come from ignoring the other type
when constructing a packing. The preceding lemmas prove
$\tc(G)\leq C_{\mathrm{dir}}$ and $\tp(G)\geq K_{\mathrm{dir}}$.
Thus a sufficient remaining obligation is
\begin{equation}\label{eq:direct-comparison}
 C_{\mathrm{dir}}\leq2K_{\mathrm{dir}}.
\end{equation}
Its universal validity is not proved here. This comparison retains both
degrees and their intersection without charging for equalization, and
needs neither the exact cover optimization nor the earlier moment bounds.

The bounded Azure job \texttt{e32bb49cebfb45fb} checked this comparison
on all $106\,020$ tuples with
\[
 4\leq k\leq18,\quad 2\leq s<t\leq k-1,\quad
 \max(0,s+t-k)\leq u\leq s-1,\quad
 1\leq p\leq h(s),\quad1\leq q\leq h(t),
\]
and on $40\,000$ seeded samples with
$k\in\{35,60,120,250,1000,10000\}$. Every comparison passed.
The same job independently minimized the cost of balanced whole-graph
cuts over both outside-type orientations in each sampled instance;
every resulting cost satisfied Lemma~\ref{lem:separable-cover}'s bound.
It also checked four prescribed examples. These are finite arithmetic
checks and a bounded falsification attempt, not a proof
of~\eqref{eq:direct-comparison} for all parameters or an independent
audit of the written proofs. A failure of this sufficient comparison,
if found later, would not itself be a counterexample to Tuza.
The worker exited successfully and returned-source hashes matched;
its sources and reports are retained under
\texttt{direct-cover/azure/runs/e32bb49cebfb45fb/}.

A separate bounded diagnostic run, \texttt{468b9e7c023348a7}, made
$606\,528$ evaluations in six degree/multiplicity regimes at
$k=35,120,6000$ and retained $36$ candidate minima, checked with exact
rational arithmetic. None failed the original integer comparison.
However, replacing the palettes and clique estimates by smooth lower
bounds is not innocuous: at
\[
 (k,s,t,u,p,q)=(35,24,34,23,11,11)
\]
the relaxed cover minus the best relaxed packing bound is $111/136>0$.
The original integer bounds instead give a cover at most $554$ and a
packing at least $287$. An explicit whole-graph cut improves that cover
to $451$. This diagnoses a loss in the simplified comparison, not a
failure of Tuza, and motivates retaining explicit overlap-sensitive
cuts in the next comparison. The diagnostic search is not a global
optimization certificate. Its worker exited successfully, returned
source hashes matched, and both resource groups were confirmed deleted.

\subsection{A continuous comparison with explicit balanced cuts}
The separable cover deliberately forgets the overlap. The following
refinement retains it and admits an exact interval check over a bounded
domain, without placing an upper bound on $k$. The final version below is certified in Lemma~\ref{lem:interval-certificate}.

Suppose $k\geq35$, $s\leq t$, $t\geq k/2$, and the multiplicities
have been compressed. Normalize
\[
 x=s/k,\quad y=t/k,\quad w=u/k,\quad
 \alpha=p/k,\quad\beta=q/k,\quad\rho=1/k,\quad\mu=\alpha+\beta.
\]
Put $c^+=(1-\rho)^2/4$ and $f(d)=\max(0,d/2-1/8,d-1/2)$.
Two normalized cover upper bounds are
\[
 D_0=c^++\alpha f(x)+\beta f(y),\qquad
 D_1=(1-\rho)/2-y(1-y)+xw-w^2.
\]
Here and below a normalized cover bound is divided by $k^2$, and a
normalized packing bound bounds $2\tp(G)/k^2$.

There are four more explicit cover bounds. The normalized sizes of the
four regions $A,B,U,O$ are $x-w,y-w,w,1-x-y+w$. Fill a part of size
$1/2$ in one of the orders
\[
 (U,A,B,O),\quad(U,B,A,O),\quad(A,U,O,B),\quad(A,O,U,B),
\]
taking the minimum of the next region's size and the unfilled space.
Let $r_A,r_B,r_U,r_O$ be the amounts taken. For the first two orders use
\[
 D_\pi=c^++\alpha x+\beta y
       -\alpha r_A-\beta r_B-(\alpha+\beta)r_U;
\]
for the last two use
\[
 D_\pi=c^++\alpha x-\alpha r_A+\beta r_B-(\alpha-\beta)r_U.
\]
These formulas simply count spokes deleted when both types prefer the
filled part, or when their preferred parts are opposite. The rounding
argument of Lemma~\ref{lem:separable-cover} makes every such fractional
balanced cut a cover bound. No optimal ordering claim is needed.
Consequently $\tc(G)/k^2\leq D:=\min(D_0,D_1,D_\pi:\pi)$.
We now strengthen $D$ by also taking the minimum with
Lemma~\ref{lem:common-cut}'s unbalanced cover:
\[
 D_{\mathrm{common}}=c^++\alpha(x-1/2)+\beta(y-1/2)
                   -\mu^2/4+\bigl((1+\mu)/2-w\bigr)_+^2.
\]
In all subsequent occurrences, $D$ includes this additional candidate.

Define the relaxed outside bounds
\begin{align*}
 I&=\alpha(x-\rho)+\beta(y-\rho)
       -\frac{\alpha\beta w(w-\rho)}{xy},\\
 J&=\frac{\alpha x(x-\rho)+\beta y(y-\rho)
                 -(\mu-y)_+w(w-\rho)}{y},\qquad X=\max(I,J),\\
 A(r,z)&=\frac{(1-\rho)(1-2\rho)/3
                    +r x(x-\rho)+z y(y-\rho)}{1+r+z},\\
 P&=\max\left\{\frac{1-2\rho-3\rho^2}{3},\ X,\
          \frac{1-2\rho}{6}+\frac{2X}{3},\
          A(\alpha,\beta),A(\alpha,0),A(0,\beta)\right\}.
\end{align*}
Then $2\tp(G)/k^2\geq P$. To justify the two relaxations, the independent
allocation expectation $r_S\binom s2+r_T\binom t2-r_Sr_T\binom u2$
is increasing in each selection rate on $[0,1]^2$. Replacing
$p/h(s),q/h(t)$ by $p/s,q/t$ therefore only decreases it. The coupled
expectation, as a function of its palette size $H$, is
\[
 \frac{p\binom s2+q\binom t2-(p+q-H)_+\binom u2}{H}.
\]
It is decreasing on both sides of $H=p+q$, since on the lower side it
equals $\binom u2+
[p(\binom s2-\binom u2)+q(\binom t2-\binom u2)]/H$.
Replacing $h(t)$ by $t$ is therefore valid. The other bounds are exactly
the earlier clique, completion, and averaging estimates, without integer
rounding.

For the remaining large-core target, it suffices to establish $P\geq D$
on the continuous domain
\begin{gather*}
 1/4\leq x\leq y\leq1,\quad y\geq1/2,\quad
 \max(0,x+y-1)\leq w\leq x,\quad 0\leq\rho\leq1/35,\\
 \rho\leq\alpha\leq x,\quad\rho\leq\beta\leq y,
 \qquad\alpha+\beta\geq1/3+\rho.
\end{gather*}
This continuous domain includes every integer instance not handled by the
preceding analytic reductions, with
$k\geq35$, using Theorem~\ref{thm:direct-regimes} and
Corollary~\ref{cor:two-half}. It is larger than the integer domain, so a
failure at a real point would not by itself refute the integer comparison.
In particular, a complete interval subdivision would prove an unbounded
statement; a partial subdivision would not.

For exact interval evaluation it is better to cancel the common core
term $c^+$ first. The four cut costs $D_\pi-c^+$ simplify, in the same
order, to
\begin{align*}
 &\alpha(x-1/2)_+
  +\beta\bigl(y-\min(w,1/2)-(1/2-x)_+\bigr)_+,\\
 &\beta(y-1/2)+\alpha\bigl(x-\min(w,1/2)\bigr),\\
 &\alpha(x-1/2)_+
  +\beta\max\bigl(y-1/2,w-(x-1/2)_+\bigr),\\
 &\beta(y-1/2)+\alpha(w-y+1/2)_+.
\end{align*}
These follow by filling the indicated regions; feasibility implies
$x-w\leq1/2$, since $x\leq y$ and $x+y-w\leq1$.
For clarity, these piecewise simplifications can be checked without
trusting the interval program. Set $a=x-w$, $b=y-w$,
$o=1-x-y+w$, and let $R=(1/2-x)_+$. In the first order,
$r_U=\min(w,1/2)$, $r_A+r_U=\min(x,1/2)$, and
$r_B=\min(b,R)$, which gives the first formula. In the second order,
$U\cup B$ already fills the half-part because $y\geq1/2$; thus
$r_A=0$ and $r_U+r_B=1/2$, giving the second formula.
For the third order, $r_A=a$ and
$r_U=w-(x-1/2)_+$. The remaining space before $O,B$ is $R$, so
$r_B=(R-o)_+$. Hence
\[
 r_U+r_B=\max\{y-1/2,\ w-(x-1/2)_+\}.
\]
Indeed, if $R=0$ then $x\geq1/2$ and $r_U=w-x+1/2\geq y-1/2$
by $x+y-w\leq1$; otherwise $r_U=w$ and $R-o=y-w-1/2$.
For the fourth order, $A\cup O$ has size $1-y\leq1/2$.
The available space for $U,B$ is $y-1/2$, so
$r_U=\min(w,y-1/2)$ and $r_U+r_B=y-1/2$; the deleted $S$ mass is
$w-r_U=(w-y+1/2)_+$. This gives the last formula.

For the shadow cover the corresponding cost is
$1/4-\rho^2/4-y(1-y)+w(x-w)$.
For either outside estimate write $X_i=X_{i,0}-\rho X_{i,1}$,
$i\in\{I,J\}$. Then its outside and completion gaps above $c^+$ are
\[
 X_{i,0}-\tfrac14+\rho(\tfrac12-X_{i,1})-\tfrac14\rho^2,
 \qquad
 -\tfrac1{12}+\tfrac23X_{i,0}
 +\rho(\tfrac16-\tfrac23X_{i,1})-\tfrac14\rho^2.
\]
The clique gap is $1/12-\rho/6-5\rho^2/4$. For averaging,
\[
 A(r,z)-c^+=
 \frac{\tfrac1{12}+r(x^2-\tfrac14)+z(y^2-\tfrac14)
 +\rho[-\tfrac12+r(\tfrac12-x)+z(\tfrac12-y)]
 +\rho^2[\tfrac5{12}-(r+z)/4]}{1+r+z}.
\]
This cancellation preserves the exact inequality and avoids charging
twice for interval uncertainty in the same core term.

\begin{lemma}[Exact interval certificate]\label{lem:interval-certificate}
On the continuous domain displayed above, $P\geq D$, with the unbalanced
common-neighborhood cover included in $D$.
\end{lemma}
\begin{proof}[Computer-assisted proof]
The checker represents an interval as $[a/2^{40},b/2^{40}]$ with integer
endpoints. Addition, multiplication, squaring and division by positive
intervals round their lower endpoints down and upper endpoints up using
integer arithmetic. The rules for minimum and maximum are their endpoint
rules. Thus all enclosing intervals are rigorous; floating-point arithmetic
is not used in any mathematical inequality.

The root intervals enclose the displayed domain. Necessary inequalities
$x\leq y$, $w\leq x$, $w\geq x+y-1$, $\rho\leq\alpha\leq x$,
$\rho\leq\beta\leq y$, and $\mu\geq1/3+\rho$ contract a box without
removing any feasible point. The non-dyadic root endpoint $1/35$ is
rounded upwards and remains below $1/34$.

For a leaf, the checker either verifies the sufficient inequality of
Corollary~\ref{cor:common-cut}, or encloses all the packing gaps and cover
costs above $c^+$. It accepts the leaf only when the maximum of the lower
packing endpoints is at least the minimum of the upper cover endpoints.
The auxiliary estimate $0\leq w(x-w)\leq x^2/4$ sharpens a cover interval.
All these checks are sufficient for $P\geq D$ on the entire feasible
part of the box. Otherwise a coordinate interval is bisected, and both
children are retained.

Four adjacent $x$ intervals partition the root domain. The completed
forests have the following numbers of accepted leaves:
\[
\begin{array}{c|r}
 x\text{ interval}&\text{accepted leaves}\\\hline
 {[1/4,7/16]}&6\,537\\
 {[7/16,5/8]}&149\,989\\
 {[5/8,13/16]}&70\,801\\
 {[13/16,1]}&837
\end{array}
\]
All $228\,164$ leaves are accepted, with no unresolved leaf. The
$456\,324$-node forests were saved and replayed: every split, domain
contraction, and terminal inequality was checked again. The completed-square
identity in Corollary~\ref{cor:common-cut} was also verified by exact
rational polynomial arithmetic. Azure job \texttt{d5719db28c9b48f4}
exited successfully and its returned sources and certificate digests
matched. The four compressed trees, checker, symbolic check and reports
are retained under
\texttt{interval-unbalanced/azure/runs/d5719db28c9b48f4/}.
The four root intervals cover the whole domain, and the finite trees
cover their roots, proving the claimed inequality. This is an exact
interval certificate for a continuous domain, not a sampling argument.
A second implementation, Azure job \texttt{53fb6bee74b8468e}, subsequently
rechecked all $456\,324$ nodes and $228\,164$ leaves. It imports none of the
search program, reconstructs contractions by generic linear inequality
projection, and evaluates the leaf bounds with $64$-bit dyadic outward
precision. Sixteen rational polynomial identities check the link between
the literal normalized bounds and their canceled forms. Truncated trees,
unresolved nodes, invalid opcodes, false exclusions, a false root inequality,
and a deliberately false polynomial identity were rejected. The separate
source and report are retained under
\texttt{certificate-audit/azure/runs/53fb6bee74b8468e/}.
This verifies the scalar certificate by a second implementation. It is not
an independent review of the graph lemmas and is not Lean verification.
\end{proof}

\begin{theorem}[Every large core with at most two types]
\label{thm:two-types-large}
Every split graph with a clique part of size $k\geq35$ and at most two
active independent-side neighborhood types satisfies $\tc(G)\leq2\tp(G)$.
\end{theorem}
\begin{proof}
Delete outside vertices of degree at most one, which participate in no
triangle, and compress multiplicities by Theorem~\ref{thm:compression}.
The clique and the one-type case follow from the preceding bounds,
including Lemma~\ref{lem:single-type-certificate}. For two types, relabel
so that $s\leq t$.
If $s\leq k/4$ or $3(p+q)\leq k+3$, apply
Theorem~\ref{thm:direct-regimes}. If $t\leq k/2$, apply
Corollary~\ref{cor:two-half}. Every other instance normalizes to a point
in the certified domain: compression ensures $p\leq s$, $q\leq t$,
and the overlap inequalities are exactly set feasibility.
The packing and cover bounds give
\[
 2\tp(G)/k^2\geq P\geq D\geq\tc(G)/k^2.
\]
No upper bound on $k$ is involved.
\end{proof}

\begin{lemma}[Finite unequal-size remainder]\label{lem:unequal-finite}
Every two-type instance with $k\leq34$ and unequal, crossing active
neighborhoods satisfies Tuza's inequality.
\end{lemma}
\begin{proof}[Computer-assisted proof]
Compress multiplicities and relabel the types so that $s<t$.
Crossing means $S\setminus T$ and $T\setminus S$ are both nonempty.
Thus all possibilities are included in the finite domain
\[
 4\leq k\leq34,\quad 2\leq s<t\leq k-1,\quad
 \max(0,s+t-k)\leq u\leq s-1,\quad
 1\leq p\leq h(s),\quad1\leq q\leq h(t).
\]
Azure job \texttt{9e59240ec1104f49} exhausted all $5\,102\,568$ tuples
in this domain. The previously proved cover and packing bounds are
evaluated by integer arithmetic: a cut gives an upper bound on $\tc$,
and clique packing, independent/coupled colors, completion, and the
three averaging bounds give lower bounds on $\tp$. For every tuple the
cover bound is at most twice the packing bound. In fact the envelope
comparison~\eqref{eq:direct-comparison} already passes every tuple;
no improved unbalanced cover is needed in this finite range.
Every $k=4,\ldots,34$ completed, the worker exited zero, and returned
source hashes matched. The checker and per-core summaries are retained
under \texttt{unequal-finite/azure/runs/9e59240ec1104f49/}.
The ordered transcript digest is
\begin{center}\small\ttfamily
2839612bb9e3957b588eb2df902635b28\\
027caa098eebdebda76f18c49b3bf14.
\end{center}
The domain exhausts all compressed unequal crossing cases, and the
bounds are valid for every graph with the given parameters. Compression
preserves both optima, so the conclusion also holds before compression.
\end{proof}

\begin{lemma}[Unified finite remainder]\label{lem:unified-finite}
Every split graph with a clique part of size at most $34$ and at most
two active outside neighborhood types satisfies $\tc(G)\leq2\tp(G)$.
\end{lemma}
\begin{proof}[Computer-assisted proof]
For $k\leq1$ there is no triangle. Otherwise compress multiplicities
using Theorem~\ref{thm:compression} and order the two neighborhood sizes.
A fresh checker exhausts
\[
 2\leq k\leq34,\quad 2\leq s\leq t\leq k,\quad
 \max(0,s+t-k)\leq u\leq s,\quad
 1\leq p\leq h(s),\quad1\leq q\leq h(t),
\]
and also all one-type cases $2\leq s\leq k$, $0\leq p\leq h(s)$.
The $p=0$ entries include the clique. The domain deliberately includes
$S=T$ and nested sets, so no balancing reduction or threshold-graph
theorem is required. Every graph in the stated class is represented;
some degenerate graphs have more than one representation, which is harmless.

For each tuple, compare twice the maximum of the proved clique, outside,
completion, and three averaging bounds with the minimum of the separable
cover bound (rounded down) and an explicitly enumerated shadow cut.
The latter enumerates eight placements of the three non-$O$ regions
and every integer split of $O$. Every candidate is a valid retained cut;
no assertion of optimality or exact triangle-cover formula is needed.
The checker imports neither of the previous finite checkers.
The integer arithmetic comparison passed throughout:
\[
\begin{array}{l|r}
 \text{parameter category}&\text{tuples checked}\\\hline
 \text{one type or clique}&7\,378\\
 \text{nested neighborhoods (including identical)}&1\,274\,881\\
 \text{equal-size crossing neighborhoods}&668\,760\\
 \text{unequal-size crossing neighborhoods}&5\,102\,568\\\hline
 \text{total}&7\,053\,587
\end{array}
\]
No tuple required the optional whole-graph cut fallback, and no tuple
failed. All $k=2,\ldots,34$ completed. As a separate check of the
implementation, direct exact optimization on $30$ small labeled graphs
confirmed that the claimed packing bounds did not exceed the optimum and
the cover bounds did not fall below it; these controls do not replace
the general graph lemmas.

Azure job \texttt{2a9f84bf04484c4e} exited zero, with source hashes verified.
The source, per-$k$ transcript hashes, counts, and report are retained in
\texttt{unified-finite/azure/runs/2a9f84bf04484c4e/}.
Compression preserves both parameters, so the exhaustive compressed
comparison proves the stated finite remainder.
\end{proof}

\section{Removing the exact clique-packing formula}
The direct sum-coloring lemma gives the explicit clique lower bound
\[
 g(n)=\begin{cases}0,&n\leq2,\\
 \left\lceil(n-1)(n-2)/6\right\rceil,&n\geq3.
 \end{cases}
\]
The value $g(n)$ is a guaranteed packing size, not a claim about the
optimum. In particular, for $k\geq3$,
\[
 2\tp(K_k)\geq2g(k)\geq\frac{k^2-3k+2}{3}.
\]
For any nonempty finite graph $H$, the same lemma gives
$\tp(H)\geq\lceil t(H)/|V(H)|\rceil$ directly.
The stronger permutation-averaging bound in Lemma~\ref{lem:averaging}
is not used in the final comparisons.

\begin{lemma}[Analytic reductions using sums of labels]
\label{lem:cyclic-reductions}
Let $G$ have at most two active neighborhood types and $k\geq43$.
Compress multiplicities, write $m=p+q$, and order $s\leq t$.
Tuza's inequality follows from the sum-coloring packing bounds if there is at
most one active type, or if any of the following holds:
\[
 s\leq k/4,\qquad 3m\leq k+5,\qquad t\leq k/2.
\]
\end{lemma}
\begin{proof}
First consider a single type of size $d$ with $r\leq h(d)$ centers.
The local-color construction, greedy completion, and direct sum-coloring
bound give
\[
 2\tp(H)\geq\max\left\{2g(k),\ r(d-1),\
 \frac{2c_k+2r(d-1)}3,\
 \frac{k(k-1)(k-2)/3+rd(d-1)}{k+r}\right\}.
\]
We prove that this maximum is at least $c_k+r f_k(d)$.
For $d\leq k/4$, the clique estimate suffices, since
\[
 2g(k)-c_k\geq\frac{k^2-6k+5}{12}>0.
\]
For $k/4\leq d\leq3k/4$ and $r\geq(k+5)/3$, combine one quarter
of the clique bound with three quarters of the completion bound. The
margin over the separable cover is at least
\[
 \frac{g(k)-c_k}{2}+\frac{r(k-4)}8
 \geq\frac{k-19}{24}>0.
\]
In the complementary range use the displayed algebraic identities for
$F_d(m)$ and $P(k,m)$ in the proof of
Lemma~\ref{lem:single-type-certificate}. The identities themselves make
no use of a clique-packing theorem. The polynomial $P$ decreases in
$m\geq0$, and its new endpoint is positive for $k\geq43$:
\begin{align*}
 9P(k,(k+5)/3)&=8k^3-342k^2+93k-785\\
 &=8h^3+690h^2+15057h+6912>0,\qquad h=k-43\geq0.
\end{align*}
Therefore $F_d(r)\geq0$ for $r\leq(k+5)/3$, proving the middle-degree
case using the whole-graph sum-coloring bound. For $d\geq3k/4$, retain the elementary high-degree
argument: the outside bound suffices when $r\geq(k+2)/2$, while the
identity for $Q(k,r)$ and its positive endpoint prove the remaining
range using that same bound. Neither of these high-degree steps uses the exact
clique formula. Thus
\[
 2\tp(H)\geq c_k+r f_k(d)
\]
holds for a single type, including zero multiplicity.

For two types with $s\leq k/4$, the separable cover is at most
$c_k+qf_k(t)$, which the just-proved single-type packing bound pays for.
If $3m\leq k+5$ and $s,t>k/4$, average $F_s(m)$ and $F_t(m)$ with
weights $p/m,q/m$. Both are nonnegative: use the new endpoint for $P$
in the middle-degree range, and $m\leq(k+5)/3\leq(k+2)/2$ for $Q$
in the high-degree range. This gives the required cover--packing
comparison exactly as in the weighted calculation of
Theorem~\ref{thm:direct-regimes}, using only the direct whole-graph bound.

Finally, if $t\leq k/2$, the explicit shadow cut in
Lemma~\ref{lem:unequal-shadow} gives $\tc(G)\leq c_k+k^2/16$, whereas
\[
 2g(k)-c_k-k^2/16\geq\frac{k^2-24k+20}{48}>0
 \qquad(k\geq43).
\]
This last comparison is another application of sum-coloring in a clique,
with no exact optimum formula.
\end{proof}

\begin{lemma}[The cyclic continuous certificate]\label{lem:cyclic-interval}
Retain the normalized cover bound $D$ and the outside, completion, and
whole-graph bounds from Lemma~\ref{lem:interval-certificate}. Replace only
the clique term in $P$ by $(1-\rho)(1-2\rho)/3$, and call the resulting
maximum $P_{\mathrm{cyc}}$. Then $P_{\mathrm{cyc}}\geq D$ on
\begin{gather*}
 1/4\leq x\leq y\leq1,\quad y\geq1/2,\quad
 \max(0,x+y-1)\leq w\leq x,\quad 0\leq\rho\leq1/43,\\
 \rho\leq\alpha\leq x,\quad \rho\leq\beta\leq y,\qquad
 3(\alpha+\beta)\geq1+5\rho.
\end{gather*}
Every packing term here follows without the exact clique-packing formula.
\end{lemma}
\begin{proof}[Computer-assisted proof]
The clique term is precisely $(k-1)(k-2)/(3k^2)$, now supplied by
Lemma~\ref{lem:sum-coloring}. The formerly averaged terms use only
$\tp(H)\geq t(H)/|V(H)|$, supplied by the same direct lemma.
Thus their numerical expressions are identical to those certified
in this section; the new argument does not weaken a continuous bound.
All outside and completion bounds are unchanged. After canceling
$c^+=(1-\rho)^2/4$, the new clique gap is
\[
 \frac1{12}-\frac\rho2+\frac{5\rho^2}{12}.
\]
The sufficient common-neighborhood criterion remains valid, because its
proof uses only $\tp(H)\geq t(H)/|V(H)|$ and $\rho\leq1/34$.

A fresh exact interval run uses the new clique gap, the root endpoint
$\rho\leq1/43$, and the new necessary inequality
$3(\alpha+\beta)\geq1+5\rho$. Endpoints are dyadic multiples of
$2^{-40}$. The new multiplicity contractions round lower endpoints down
and upper endpoints up when dividing by $3$ or $5$. All other enclosing,
bisection, and acceptance rules are as described in the earlier
certificate proof. The resulting adjacent slabs are fully certified:
\[
\begin{array}{c|r}
 x\text{ interval}&\text{accepted leaves}\\\hline
 {[1/4,7/16]}&5\,604\\
 {[7/16,5/8]}&116\,226\\
 {[5/8,13/16]}&55\,861\\
 {[13/16,1]}&815\\\hline
 \text{total}&178\,506
\end{array}
\]
All $357\,008$ nodes were replayed, with no unresolved leaf.
The new clique-gap identity, the revised endpoint polynomial at $43$,
and the two revised analytic margin identities were checked by exact
rational polynomial arithmetic. Thus the calculation certifies the
entire continuous domain and places no upper bound on $k$.

Azure job \texttt{17ed3df4975c4594} returned these complete trees with
matching source and certificate hashes. Sources and reports are retained
under \texttt{cyclic-proof/azure/runs/17ed3df4975c4594/}.
The earlier certificate is not used to certify the changed inequalities.
A separate implementation in Azure job \texttt{93cea9de41fc4180}
rechecked every new node and leaf with $64$-bit outward dyadic intervals
and generic linear constraint projection. All passed. It also checked
$16$ exact polynomial identities and rejected deliberately invalid trees
and a false polynomial identity. This second implementation verifies the
scalar certificate; it is not a Lean proof or an independent review of
the graph-theoretic reductions.
\end{proof}

\begin{lemma}[The finite certificate using sum-coloring]\label{lem:cyclic-finite}
Every split graph with at most two active neighborhood types and a clique
part of size $k\leq42$ satisfies Tuza's inequality using sum-coloring
and the elementary graph bounds in these notes.
\end{lemma}
\begin{proof}[Computer-assisted proof]
For $k\leq1$ there is no triangle. Compress multiplicities and, for two
types, order $s\leq t$. The finite domain is
\[
 2\leq k\leq42,\quad 2\leq s\leq t\leq k,\quad
 \max(0,s+t-k)\leq u\leq s,\quad
 1\leq p\leq h(s),\quad1\leq q\leq h(t),
\]
with the separate one-type cases $2\leq s\leq k$, $0\leq p\leq h(s)$.
The clique itself contributes $g(k)$. For each of the three induced
graphs with $(r,z)=(p,q),(p,0),(0,q)$, use the direct bound
\[
 \left\lceil\frac{\binom k3+r\binom s2+z\binom t2}{k+r+z}\right\rceil.
\]
This may be smaller than the earlier permutation-averaging bound, so
the entire finite domain was checked afresh. The covers, local-color
bounds, and greedy completion remain valid as before.
The complete integer comparison passed in every category:
\[
\begin{array}{l|r}
 \text{one type or clique}&13\,622\\
 \text{nested neighborhoods (including identical)}&3\,591\,105\\
 \text{equal-size crossing neighborhoods}&1\,920\,614\\
 \text{unequal-size crossing neighborhoods}&18\,323\,030\\\hline
 \text{total parameter tuples}&23\,848\,371
\end{array}
\]
These are parameter representations, not a count of nonisomorphic graphs.
In $60$ tuples the separable and shadow envelope alone was insufficient;
an explicit whole-graph cut supplies the required cover. To construct
that cut, assign the two outside types either the same side or opposite
sides, order the core vertices by their spoke weights, and inspect every
prefix length $j=0,\ldots,k$. Every prefix is a genuine bipartition;
its deleted-edge count is
\[
 \binom k2-j(k-j)+p(s-\ell_S)+
 \begin{cases}q(t-\ell_T),&\text{same outside side},\\
 q\ell_T,&\text{opposite outside sides},\end{cases}
\]
where $\ell_S,\ell_T$ count neighbors in the prefix. Optimality of the
ordering is unnecessary: any candidate whose cost is at most twice the
packing lower bound suffices. No exact cover formula is used.

Every $k=2,\ldots,42$ completed with no failed tuple in Azure job
\texttt{86a6063f11ee4e5a}; the scalar scan took $31.37$ seconds. Its
source, per-$k$ transcript digests, summaries, and all $60$ additional
cut cases are retained under \texttt{lean-sum-coloring/azure/runs/}.
For each extra case, a separate labeled-graph implementation saved
the core partition and outside orientations, counted the deleted edges
directly, and verified that every triangle met the cover. These explicit
witnesses are retained with the new job as
\texttt{fallback-covers.json}. The scalar source selects the $60$ cases;
the labeled-graph check validates their cuts
independently of its cost formula. This is not a reimplementation of all
$23\,848\,371$ finite scalar comparisons. The earlier cyclic-bound run
\texttt{17ed3df4975c4594}, its $38$ additional cuts, and their separate
audit \texttt{93cea9de41fc4180} remain historical evidence for the previous
stronger finite bounds; the present finite claim uses the new run.
\end{proof}

\begin{theorem}[Large cores from elementary packings]
\label{thm:cyclic-large}
Every split graph with at most two active neighborhood types and
$k\geq43$ satisfies $\tc(G)\leq2\tp(G)$ without using the exact
clique-packing formula.
\end{theorem}
\begin{proof}
Compress the multiplicities. Apply Lemma~\ref{lem:cyclic-reductions} to
its stated cases. Every remaining graph, after ordering $s\leq t$,
normalizes to a point in the domain of Lemma~\ref{lem:cyclic-interval}.
All normalized bounds follow from the explicit constructions, and
\[
 2\tp(G)/k^2\geq P_{\mathrm{cyc}}\geq D\geq\tc(G)/k^2.
\]
This proves the result for every $k\geq43$ simultaneously.
\end{proof}

\begin{theorem}[At most two active neighborhood types]
\label{thm:two-types-all}
Let $G$ be a finite simple graph with a split partition
$V(G)=C\mathbin{\dot\cup}I$. Suppose the vertices of $I$ having at least
two neighbors in $C$ realize at most two distinct neighborhoods in $C$.
Then
\[
 \boxed{\tc(G)\leq2\tp(G).}
\]
There is no restriction on $|C|$, the two neighborhood sizes, their
intersection, or their multiplicities.
\end{theorem}
\begin{proof}
Delete the outside vertices of degree at most one. For $k=|C|\leq42$,
apply Lemma~\ref{lem:cyclic-finite}; for $k\geq43$, apply
Theorem~\ref{thm:cyclic-large}. These alternatives exhaust the stated
class. Both branches obtain their clique and whole-graph packing bounds
from Lemma~\ref{lem:sum-coloring}, without an exact maximum-packing
formula or permutation averaging.
\end{proof}

\section{Formal verification}
Theorem~\ref{thm:two-types-all} is now fully formalized in Lean 4.34.0,
with Mathlib pinned to revision \texttt{5ed296525643}.
The public theorem is
\texttt{TuzaTwoTypes.tuza\_of\_two\_active\_types} in
\path{lean-finish-small/Result.lean}. Its only mathematical hypotheses specify
a finite simple graph, a split partition, and at most two distinct active
outside neighborhoods. There is no bound on the clique size, overlap,
or multiplicities. The conclusion uses the actual optimal triangle
edge-cover and triangle-packing numbers.

The formal graph proof includes the packing and cover constructions,
compression of both parameters, graph representation, removal of inactive
vertices, rounding, all scalar bounds and the normalized-domain reductions.
Its final route uses matching decompositions, explicit cuts, Mantel's
bound through Mathlib's Tur\'an theorem, and sum-coloring of triangles.
The exact clique-packing formula, threshold-graph theorem, exact cover
formula and permutation averaging are not inputs to this route.

The finite certificate covers all $23\,848\,371$ required compressed
parameter tuples for $2\leq k\leq42$. Its $2326$ checked modules include
the range-assembly proofs and the quantified final lemma. The continuous
certificate has $178\,506$ leaves, divided into $347$ subtree modules,
four root modules and one domain-coverage module, with additional helper
modules bounding compiler memory. It proves the required
inequality throughout the normalized real domain and therefore covers
every $k\geq43$ simultaneously. The small-core cases are handled directly.

Python generates candidate proof terms and transports certificate files;
its output is not an axiom. Lean checks the arithmetic proofs with
\texttt{decide +kernel}, as well as the interval soundness, range joins,
domain coverage and final graph theorem. The project proof sources contain
no \texttt{sorry}, \texttt{admit}, custom axioms, \texttt{unsafe},
\texttt{native\_decide} or \texttt{implemented\_by}. The final theorem's
reported axiom dependencies are exactly \texttt{propext},
\texttt{Classical.choice} and \texttt{Quot.sound}.
A deliberately false arithmetic proof was rejected.

All compilation and certificate replay took place in bounded Azure jobs.
The final verification report is retained at
\path{lean-finish-small/azure/runs/fa2b8345ae1d465f/verification-report.json},
with source hashes, compiled-object hashes, compiler logs and a separate
resource-deletion report. Previously checked modules were reused only
with matching source and object hashes and the same toolchain pins.
This establishes the stated two-type subclass theorem; it does not resolve
Tuza's conjecture for arbitrary graphs. No priority claim is made.

The earlier sections also record the development of stronger bounds and
exploratory computations. The following checks are historical context;
the formal theorem above no longer depends on their Python success flags.

\paragraph{Earlier exploratory checks.}
The companion scripts check a fixed list of small witnesses, their cover
counts, the allocation-bound arithmetic, and the exact polynomial identities
in Theorem~\ref{thm:few-outside}. The new fixed checker verifies a $42$-triangle
packing and a $75$-edge cover, and the moment formulas for two tiny color
experiments (six and sixty assignments). These are bounded illustrations
and negative controls, not computational verification of the general
theorems. Two bounded Azure scalar searches were run: the first found the
old-bound failure after $54\,646$ evaluations; the second found no failure
of the strengthened comparison in $678\,730$ evaluations, consisting of a
complete specified sweep for $9\leq k\leq24$ and $96$ seeded larger cases.
The sweep uses crossing types, symmetry and the proved multiplicity caps;
its exact domain is specified by the saved worker source. These searches
do not verify all clique sizes or the written proofs. The fifth-gate checker
adds two tiny coupled-color experiments (six and three assignments) and
a fixed $61$-triangle packing with a $75$-edge cover. The fixed greedy
completion and its triangle-free residue are checked directly. At that stage, no new
Azure job or family enumeration was needed, and no Lean build was run.
The sixth gate used a separate bounded Azure job for the finite remainder
$k\leq34$ and six exact polynomial identities, as described above. This
is an exhaustive remainder after a uniform theorem, unlike the earlier
exploratory searches. Its separate runtime and cleanup reports are retained.
The missing-edge allowance uses the established clique packing bound
explicitly cited above; that design-theoretic ingredient is not proved
afresh in these notes.

\begin{thebibliography}{9}
\bibitem{bonamy} M. Bonamy et al., \emph{Tuza's Conjecture for Threshold
Graphs}, Discrete Mathematics \& Theoretical Computer Science 24(1), 2022.
\url{https://arxiv.org/abs/2105.09871}.
\bibitem{zeng} Z. Zeng, \emph{Tuza's Conjecture for Split Graphs with an
Eight-Vertex Clique Part and Two Neighborhood Types}, preprint, 2026.
\url{https://www.preprints.org/manuscript/202608.1304}.
\bibitem{dense} L. Chahua and J. Guti\'errez,
\emph{On Tuza's conjecture in dense graphs},
Discrete Applied Mathematics 377 (2025), 225--233.
\url{https://arxiv.org/abs/2405.11409}.
\bibitem{mathlib-turan} Mathlib contributors, \emph{Tur\'an's theorem},
\path{Mathlib/Combinatorics/SimpleGraph/Extremal/Turan.lean},
pinned revision \texttt{5ed296525643}.
\url{https://github.com/leanprover-community/mathlib4/blob/5ed2965256430c3649e86755f9576b54eca72435/Mathlib/Combinatorics/SimpleGraph/Extremal/Turan.lean}.
\end{thebibliography}
\end{document}
